\documentclass[11pt]{article}

    \usepackage[breakable]{tcolorbox}
    \usepackage{parskip} % Stop auto-indenting (to mimic markdown behaviour)
    

    % Basic figure setup, for now with no caption control since it's done
    % automatically by Pandoc (which extracts ![](path) syntax from Markdown).
    \usepackage{graphicx}
    % Keep aspect ratio if custom image width or height is specified
    \setkeys{Gin}{keepaspectratio}
    % Maintain compatibility with old templates. Remove in nbconvert 6.0
    \let\Oldincludegraphics\includegraphics
    % Ensure that by default, figures have no caption (until we provide a
    % proper Figure object with a Caption API and a way to capture that
    % in the conversion process - todo).
    \usepackage{caption}
    \DeclareCaptionFormat{nocaption}{}
    \captionsetup{format=nocaption,aboveskip=0pt,belowskip=0pt}

    \usepackage{float}
    \floatplacement{figure}{H} % forces figures to be placed at the correct location
    \usepackage{xcolor} % Allow colors to be defined
    \usepackage{enumerate} % Needed for markdown enumerations to work
    \usepackage{geometry} % Used to adjust the document margins
    \usepackage{amsmath} % Equations
    \usepackage{amssymb} % Equations
    \usepackage{textcomp} % defines textquotesingle
    % Hack from http://tex.stackexchange.com/a/47451/13684:
    \AtBeginDocument{%
        \def\PYZsq{\textquotesingle}% Upright quotes in Pygmentized code
    }
    \usepackage{upquote} % Upright quotes for verbatim code
    \usepackage{eurosym} % defines \euro

    \usepackage{iftex}
    \ifPDFTeX
        \usepackage[T1]{fontenc}
        \IfFileExists{alphabeta.sty}{
              \usepackage{alphabeta}
          }{
              \usepackage[mathletters]{ucs}
              \usepackage[utf8x]{inputenc}
          }
    \else
        \usepackage{fontspec}
        \usepackage{unicode-math}
    \fi

    \usepackage{fancyvrb} % verbatim replacement that allows latex
    \usepackage{grffile} % extends the file name processing of package graphics
                         % to support a larger range
    \makeatletter % fix for old versions of grffile with XeLaTeX
    \@ifpackagelater{grffile}{2019/11/01}
    {
      % Do nothing on new versions
    }
    {
      \def\Gread@@xetex#1{%
        \IfFileExists{"\Gin@base".bb}%
        {\Gread@eps{\Gin@base.bb}}%
        {\Gread@@xetex@aux#1}%
      }
    }
    \makeatother
    \usepackage[Export]{adjustbox} % Used to constrain images to a maximum size
    \adjustboxset{max size={0.9\linewidth}{0.9\paperheight}}

    % The hyperref package gives us a pdf with properly built
    % internal navigation ('pdf bookmarks' for the table of contents,
    % internal cross-reference links, web links for URLs, etc.)
    \usepackage{hyperref}
    % The default LaTeX title has an obnoxious amount of whitespace. By default,
    % titling removes some of it. It also provides customization options.
    \usepackage{titling}
    \usepackage{longtable} % longtable support required by pandoc >1.10
    \usepackage{booktabs}  % table support for pandoc > 1.12.2
    \usepackage{array}     % table support for pandoc >= 2.11.3
    \usepackage{calc}      % table minipage width calculation for pandoc >= 2.11.1
    \usepackage[inline]{enumitem} % IRkernel/repr support (it uses the enumerate* environment)
    \usepackage[normalem]{ulem} % ulem is needed to support strikethroughs (\sout)
                                % normalem makes italics be italics, not underlines
    \usepackage{soul}      % strikethrough (\st) support for pandoc >= 3.0.0
    \usepackage{mathrsfs}
    

    
    % Colors for the hyperref package
    \definecolor{urlcolor}{rgb}{0,.145,.698}
    \definecolor{linkcolor}{rgb}{.71,0.21,0.01}
    \definecolor{citecolor}{rgb}{.12,.54,.11}

    % ANSI colors
    \definecolor{ansi-black}{HTML}{3E424D}
    \definecolor{ansi-black-intense}{HTML}{282C36}
    \definecolor{ansi-red}{HTML}{E75C58}
    \definecolor{ansi-red-intense}{HTML}{B22B31}
    \definecolor{ansi-green}{HTML}{00A250}
    \definecolor{ansi-green-intense}{HTML}{007427}
    \definecolor{ansi-yellow}{HTML}{DDB62B}
    \definecolor{ansi-yellow-intense}{HTML}{B27D12}
    \definecolor{ansi-blue}{HTML}{208FFB}
    \definecolor{ansi-blue-intense}{HTML}{0065CA}
    \definecolor{ansi-magenta}{HTML}{D160C4}
    \definecolor{ansi-magenta-intense}{HTML}{A03196}
    \definecolor{ansi-cyan}{HTML}{60C6C8}
    \definecolor{ansi-cyan-intense}{HTML}{258F8F}
    \definecolor{ansi-white}{HTML}{C5C1B4}
    \definecolor{ansi-white-intense}{HTML}{A1A6B2}
    \definecolor{ansi-default-inverse-fg}{HTML}{FFFFFF}
    \definecolor{ansi-default-inverse-bg}{HTML}{000000}

    % common color for the border for error outputs.
    \definecolor{outerrorbackground}{HTML}{FFDFDF}

    % commands and environments needed by pandoc snippets
    % extracted from the output of `pandoc -s`
    \providecommand{\tightlist}{%
      \setlength{\itemsep}{0pt}\setlength{\parskip}{0pt}}
    \DefineVerbatimEnvironment{Highlighting}{Verbatim}{commandchars=\\\{\}}
    % Add ',fontsize=\small' for more characters per line
    \newenvironment{Shaded}{}{}
    \newcommand{\KeywordTok}[1]{\textcolor[rgb]{0.00,0.44,0.13}{\textbf{{#1}}}}
    \newcommand{\DataTypeTok}[1]{\textcolor[rgb]{0.56,0.13,0.00}{{#1}}}
    \newcommand{\DecValTok}[1]{\textcolor[rgb]{0.25,0.63,0.44}{{#1}}}
    \newcommand{\BaseNTok}[1]{\textcolor[rgb]{0.25,0.63,0.44}{{#1}}}
    \newcommand{\FloatTok}[1]{\textcolor[rgb]{0.25,0.63,0.44}{{#1}}}
    \newcommand{\CharTok}[1]{\textcolor[rgb]{0.25,0.44,0.63}{{#1}}}
    \newcommand{\StringTok}[1]{\textcolor[rgb]{0.25,0.44,0.63}{{#1}}}
    \newcommand{\CommentTok}[1]{\textcolor[rgb]{0.38,0.63,0.69}{\textit{{#1}}}}
    \newcommand{\OtherTok}[1]{\textcolor[rgb]{0.00,0.44,0.13}{{#1}}}
    \newcommand{\AlertTok}[1]{\textcolor[rgb]{1.00,0.00,0.00}{\textbf{{#1}}}}
    \newcommand{\FunctionTok}[1]{\textcolor[rgb]{0.02,0.16,0.49}{{#1}}}
    \newcommand{\RegionMarkerTok}[1]{{#1}}
    \newcommand{\ErrorTok}[1]{\textcolor[rgb]{1.00,0.00,0.00}{\textbf{{#1}}}}
    \newcommand{\NormalTok}[1]{{#1}}

    % Additional commands for more recent versions of Pandoc
    \newcommand{\ConstantTok}[1]{\textcolor[rgb]{0.53,0.00,0.00}{{#1}}}
    \newcommand{\SpecialCharTok}[1]{\textcolor[rgb]{0.25,0.44,0.63}{{#1}}}
    \newcommand{\VerbatimStringTok}[1]{\textcolor[rgb]{0.25,0.44,0.63}{{#1}}}
    \newcommand{\SpecialStringTok}[1]{\textcolor[rgb]{0.73,0.40,0.53}{{#1}}}
    \newcommand{\ImportTok}[1]{{#1}}
    \newcommand{\DocumentationTok}[1]{\textcolor[rgb]{0.73,0.13,0.13}{\textit{{#1}}}}
    \newcommand{\AnnotationTok}[1]{\textcolor[rgb]{0.38,0.63,0.69}{\textbf{\textit{{#1}}}}}
    \newcommand{\CommentVarTok}[1]{\textcolor[rgb]{0.38,0.63,0.69}{\textbf{\textit{{#1}}}}}
    \newcommand{\VariableTok}[1]{\textcolor[rgb]{0.10,0.09,0.49}{{#1}}}
    \newcommand{\ControlFlowTok}[1]{\textcolor[rgb]{0.00,0.44,0.13}{\textbf{{#1}}}}
    \newcommand{\OperatorTok}[1]{\textcolor[rgb]{0.40,0.40,0.40}{{#1}}}
    \newcommand{\BuiltInTok}[1]{{#1}}
    \newcommand{\ExtensionTok}[1]{{#1}}
    \newcommand{\PreprocessorTok}[1]{\textcolor[rgb]{0.74,0.48,0.00}{{#1}}}
    \newcommand{\AttributeTok}[1]{\textcolor[rgb]{0.49,0.56,0.16}{{#1}}}
    \newcommand{\InformationTok}[1]{\textcolor[rgb]{0.38,0.63,0.69}{\textbf{\textit{{#1}}}}}
    \newcommand{\WarningTok}[1]{\textcolor[rgb]{0.38,0.63,0.69}{\textbf{\textit{{#1}}}}}
    \makeatletter
    \newsavebox\pandoc@box
    \newcommand*\pandocbounded[1]{%
      \sbox\pandoc@box{#1}%
      % scaling factors for width and height
      \Gscale@div\@tempa\textheight{\dimexpr\ht\pandoc@box+\dp\pandoc@box\relax}%
      \Gscale@div\@tempb\linewidth{\wd\pandoc@box}%
      % select the smaller of both
      \ifdim\@tempb\p@<\@tempa\p@
        \let\@tempa\@tempb
      \fi
      % scaling accordingly (\@tempa < 1)
      \ifdim\@tempa\p@<\p@
        \scalebox{\@tempa}{\usebox\pandoc@box}%
      % scaling not needed, use as it is
      \else
        \usebox{\pandoc@box}%
      \fi
    }
    \makeatother

    % Define a nice break command that doesn't care if a line doesn't already
    % exist.
    \def\br{\hspace*{\fill} \\* }
    % Math Jax compatibility definitions
    \def\gt{>}
    \def\lt{<}
    \let\Oldtex\TeX
    \let\Oldlatex\LaTeX
    \renewcommand{\TeX}{\textrm{\Oldtex}}
    \renewcommand{\LaTeX}{\textrm{\Oldlatex}}
    % Document parameters
    % Document title
    \title{Ass1\_1}
    
    
    
    
    
    
    
% Pygments definitions
\makeatletter
\def\PY@reset{\let\PY@it=\relax \let\PY@bf=\relax%
    \let\PY@ul=\relax \let\PY@tc=\relax%
    \let\PY@bc=\relax \let\PY@ff=\relax}
\def\PY@tok#1{\csname PY@tok@#1\endcsname}
\def\PY@toks#1+{\ifx\relax#1\empty\else%
    \PY@tok{#1}\expandafter\PY@toks\fi}
\def\PY@do#1{\PY@bc{\PY@tc{\PY@ul{%
    \PY@it{\PY@bf{\PY@ff{#1}}}}}}}
\def\PY#1#2{\PY@reset\PY@toks#1+\relax+\PY@do{#2}}

\@namedef{PY@tok@w}{\def\PY@tc##1{\textcolor[rgb]{0.73,0.73,0.73}{##1}}}
\@namedef{PY@tok@c}{\let\PY@it=\textit\def\PY@tc##1{\textcolor[rgb]{0.24,0.48,0.48}{##1}}}
\@namedef{PY@tok@cp}{\def\PY@tc##1{\textcolor[rgb]{0.61,0.40,0.00}{##1}}}
\@namedef{PY@tok@k}{\let\PY@bf=\textbf\def\PY@tc##1{\textcolor[rgb]{0.00,0.50,0.00}{##1}}}
\@namedef{PY@tok@kp}{\def\PY@tc##1{\textcolor[rgb]{0.00,0.50,0.00}{##1}}}
\@namedef{PY@tok@kt}{\def\PY@tc##1{\textcolor[rgb]{0.69,0.00,0.25}{##1}}}
\@namedef{PY@tok@o}{\def\PY@tc##1{\textcolor[rgb]{0.40,0.40,0.40}{##1}}}
\@namedef{PY@tok@ow}{\let\PY@bf=\textbf\def\PY@tc##1{\textcolor[rgb]{0.67,0.13,1.00}{##1}}}
\@namedef{PY@tok@nb}{\def\PY@tc##1{\textcolor[rgb]{0.00,0.50,0.00}{##1}}}
\@namedef{PY@tok@nf}{\def\PY@tc##1{\textcolor[rgb]{0.00,0.00,1.00}{##1}}}
\@namedef{PY@tok@nc}{\let\PY@bf=\textbf\def\PY@tc##1{\textcolor[rgb]{0.00,0.00,1.00}{##1}}}
\@namedef{PY@tok@nn}{\let\PY@bf=\textbf\def\PY@tc##1{\textcolor[rgb]{0.00,0.00,1.00}{##1}}}
\@namedef{PY@tok@ne}{\let\PY@bf=\textbf\def\PY@tc##1{\textcolor[rgb]{0.80,0.25,0.22}{##1}}}
\@namedef{PY@tok@nv}{\def\PY@tc##1{\textcolor[rgb]{0.10,0.09,0.49}{##1}}}
\@namedef{PY@tok@no}{\def\PY@tc##1{\textcolor[rgb]{0.53,0.00,0.00}{##1}}}
\@namedef{PY@tok@nl}{\def\PY@tc##1{\textcolor[rgb]{0.46,0.46,0.00}{##1}}}
\@namedef{PY@tok@ni}{\let\PY@bf=\textbf\def\PY@tc##1{\textcolor[rgb]{0.44,0.44,0.44}{##1}}}
\@namedef{PY@tok@na}{\def\PY@tc##1{\textcolor[rgb]{0.41,0.47,0.13}{##1}}}
\@namedef{PY@tok@nt}{\let\PY@bf=\textbf\def\PY@tc##1{\textcolor[rgb]{0.00,0.50,0.00}{##1}}}
\@namedef{PY@tok@nd}{\def\PY@tc##1{\textcolor[rgb]{0.67,0.13,1.00}{##1}}}
\@namedef{PY@tok@s}{\def\PY@tc##1{\textcolor[rgb]{0.73,0.13,0.13}{##1}}}
\@namedef{PY@tok@sd}{\let\PY@it=\textit\def\PY@tc##1{\textcolor[rgb]{0.73,0.13,0.13}{##1}}}
\@namedef{PY@tok@si}{\let\PY@bf=\textbf\def\PY@tc##1{\textcolor[rgb]{0.64,0.35,0.47}{##1}}}
\@namedef{PY@tok@se}{\let\PY@bf=\textbf\def\PY@tc##1{\textcolor[rgb]{0.67,0.36,0.12}{##1}}}
\@namedef{PY@tok@sr}{\def\PY@tc##1{\textcolor[rgb]{0.64,0.35,0.47}{##1}}}
\@namedef{PY@tok@ss}{\def\PY@tc##1{\textcolor[rgb]{0.10,0.09,0.49}{##1}}}
\@namedef{PY@tok@sx}{\def\PY@tc##1{\textcolor[rgb]{0.00,0.50,0.00}{##1}}}
\@namedef{PY@tok@m}{\def\PY@tc##1{\textcolor[rgb]{0.40,0.40,0.40}{##1}}}
\@namedef{PY@tok@gh}{\let\PY@bf=\textbf\def\PY@tc##1{\textcolor[rgb]{0.00,0.00,0.50}{##1}}}
\@namedef{PY@tok@gu}{\let\PY@bf=\textbf\def\PY@tc##1{\textcolor[rgb]{0.50,0.00,0.50}{##1}}}
\@namedef{PY@tok@gd}{\def\PY@tc##1{\textcolor[rgb]{0.63,0.00,0.00}{##1}}}
\@namedef{PY@tok@gi}{\def\PY@tc##1{\textcolor[rgb]{0.00,0.52,0.00}{##1}}}
\@namedef{PY@tok@gr}{\def\PY@tc##1{\textcolor[rgb]{0.89,0.00,0.00}{##1}}}
\@namedef{PY@tok@ge}{\let\PY@it=\textit}
\@namedef{PY@tok@gs}{\let\PY@bf=\textbf}
\@namedef{PY@tok@ges}{\let\PY@bf=\textbf\let\PY@it=\textit}
\@namedef{PY@tok@gp}{\let\PY@bf=\textbf\def\PY@tc##1{\textcolor[rgb]{0.00,0.00,0.50}{##1}}}
\@namedef{PY@tok@go}{\def\PY@tc##1{\textcolor[rgb]{0.44,0.44,0.44}{##1}}}
\@namedef{PY@tok@gt}{\def\PY@tc##1{\textcolor[rgb]{0.00,0.27,0.87}{##1}}}
\@namedef{PY@tok@err}{\def\PY@bc##1{{\setlength{\fboxsep}{\string -\fboxrule}\fcolorbox[rgb]{1.00,0.00,0.00}{1,1,1}{\strut ##1}}}}
\@namedef{PY@tok@kc}{\let\PY@bf=\textbf\def\PY@tc##1{\textcolor[rgb]{0.00,0.50,0.00}{##1}}}
\@namedef{PY@tok@kd}{\let\PY@bf=\textbf\def\PY@tc##1{\textcolor[rgb]{0.00,0.50,0.00}{##1}}}
\@namedef{PY@tok@kn}{\let\PY@bf=\textbf\def\PY@tc##1{\textcolor[rgb]{0.00,0.50,0.00}{##1}}}
\@namedef{PY@tok@kr}{\let\PY@bf=\textbf\def\PY@tc##1{\textcolor[rgb]{0.00,0.50,0.00}{##1}}}
\@namedef{PY@tok@bp}{\def\PY@tc##1{\textcolor[rgb]{0.00,0.50,0.00}{##1}}}
\@namedef{PY@tok@fm}{\def\PY@tc##1{\textcolor[rgb]{0.00,0.00,1.00}{##1}}}
\@namedef{PY@tok@vc}{\def\PY@tc##1{\textcolor[rgb]{0.10,0.09,0.49}{##1}}}
\@namedef{PY@tok@vg}{\def\PY@tc##1{\textcolor[rgb]{0.10,0.09,0.49}{##1}}}
\@namedef{PY@tok@vi}{\def\PY@tc##1{\textcolor[rgb]{0.10,0.09,0.49}{##1}}}
\@namedef{PY@tok@vm}{\def\PY@tc##1{\textcolor[rgb]{0.10,0.09,0.49}{##1}}}
\@namedef{PY@tok@sa}{\def\PY@tc##1{\textcolor[rgb]{0.73,0.13,0.13}{##1}}}
\@namedef{PY@tok@sb}{\def\PY@tc##1{\textcolor[rgb]{0.73,0.13,0.13}{##1}}}
\@namedef{PY@tok@sc}{\def\PY@tc##1{\textcolor[rgb]{0.73,0.13,0.13}{##1}}}
\@namedef{PY@tok@dl}{\def\PY@tc##1{\textcolor[rgb]{0.73,0.13,0.13}{##1}}}
\@namedef{PY@tok@s2}{\def\PY@tc##1{\textcolor[rgb]{0.73,0.13,0.13}{##1}}}
\@namedef{PY@tok@sh}{\def\PY@tc##1{\textcolor[rgb]{0.73,0.13,0.13}{##1}}}
\@namedef{PY@tok@s1}{\def\PY@tc##1{\textcolor[rgb]{0.73,0.13,0.13}{##1}}}
\@namedef{PY@tok@mb}{\def\PY@tc##1{\textcolor[rgb]{0.40,0.40,0.40}{##1}}}
\@namedef{PY@tok@mf}{\def\PY@tc##1{\textcolor[rgb]{0.40,0.40,0.40}{##1}}}
\@namedef{PY@tok@mh}{\def\PY@tc##1{\textcolor[rgb]{0.40,0.40,0.40}{##1}}}
\@namedef{PY@tok@mi}{\def\PY@tc##1{\textcolor[rgb]{0.40,0.40,0.40}{##1}}}
\@namedef{PY@tok@il}{\def\PY@tc##1{\textcolor[rgb]{0.40,0.40,0.40}{##1}}}
\@namedef{PY@tok@mo}{\def\PY@tc##1{\textcolor[rgb]{0.40,0.40,0.40}{##1}}}
\@namedef{PY@tok@ch}{\let\PY@it=\textit\def\PY@tc##1{\textcolor[rgb]{0.24,0.48,0.48}{##1}}}
\@namedef{PY@tok@cm}{\let\PY@it=\textit\def\PY@tc##1{\textcolor[rgb]{0.24,0.48,0.48}{##1}}}
\@namedef{PY@tok@cpf}{\let\PY@it=\textit\def\PY@tc##1{\textcolor[rgb]{0.24,0.48,0.48}{##1}}}
\@namedef{PY@tok@c1}{\let\PY@it=\textit\def\PY@tc##1{\textcolor[rgb]{0.24,0.48,0.48}{##1}}}
\@namedef{PY@tok@cs}{\let\PY@it=\textit\def\PY@tc##1{\textcolor[rgb]{0.24,0.48,0.48}{##1}}}

\def\PYZbs{\char`\\}
\def\PYZus{\char`\_}
\def\PYZob{\char`\{}
\def\PYZcb{\char`\}}
\def\PYZca{\char`\^}
\def\PYZam{\char`\&}
\def\PYZlt{\char`\<}
\def\PYZgt{\char`\>}
\def\PYZsh{\char`\#}
\def\PYZpc{\char`\%}
\def\PYZdl{\char`\$}
\def\PYZhy{\char`\-}
\def\PYZsq{\char`\'}
\def\PYZdq{\char`\"}
\def\PYZti{\char`\~}
% for compatibility with earlier versions
\def\PYZat{@}
\def\PYZlb{[}
\def\PYZrb{]}
\makeatother


    % For linebreaks inside Verbatim environment from package fancyvrb.
    \makeatletter
        \newbox\Wrappedcontinuationbox
        \newbox\Wrappedvisiblespacebox
        \newcommand*\Wrappedvisiblespace {\textcolor{red}{\textvisiblespace}}
        \newcommand*\Wrappedcontinuationsymbol {\textcolor{red}{\llap{\tiny$\m@th\hookrightarrow$}}}
        \newcommand*\Wrappedcontinuationindent {3ex }
        \newcommand*\Wrappedafterbreak {\kern\Wrappedcontinuationindent\copy\Wrappedcontinuationbox}
        % Take advantage of the already applied Pygments mark-up to insert
        % potential linebreaks for TeX processing.
        %        {, <, #, %, $, ' and ": go to next line.
        %        _, }, ^, &, >, - and ~: stay at end of broken line.
        % Use of \textquotesingle for straight quote.
        \newcommand*\Wrappedbreaksatspecials {%
            \def\PYGZus{\discretionary{\char`\_}{\Wrappedafterbreak}{\char`\_}}%
            \def\PYGZob{\discretionary{}{\Wrappedafterbreak\char`\{}{\char`\{}}%
            \def\PYGZcb{\discretionary{\char`\}}{\Wrappedafterbreak}{\char`\}}}%
            \def\PYGZca{\discretionary{\char`\^}{\Wrappedafterbreak}{\char`\^}}%
            \def\PYGZam{\discretionary{\char`\&}{\Wrappedafterbreak}{\char`\&}}%
            \def\PYGZlt{\discretionary{}{\Wrappedafterbreak\char`\<}{\char`\<}}%
            \def\PYGZgt{\discretionary{\char`\>}{\Wrappedafterbreak}{\char`\>}}%
            \def\PYGZsh{\discretionary{}{\Wrappedafterbreak\char`\#}{\char`\#}}%
            \def\PYGZpc{\discretionary{}{\Wrappedafterbreak\char`\%}{\char`\%}}%
            \def\PYGZdl{\discretionary{}{\Wrappedafterbreak\char`\$}{\char`\$}}%
            \def\PYGZhy{\discretionary{\char`\-}{\Wrappedafterbreak}{\char`\-}}%
            \def\PYGZsq{\discretionary{}{\Wrappedafterbreak\textquotesingle}{\textquotesingle}}%
            \def\PYGZdq{\discretionary{}{\Wrappedafterbreak\char`\"}{\char`\"}}%
            \def\PYGZti{\discretionary{\char`\~}{\Wrappedafterbreak}{\char`\~}}%
        }
        % Some characters . , ; ? ! / are not pygmentized.
        % This macro makes them "active" and they will insert potential linebreaks
        \newcommand*\Wrappedbreaksatpunct {%
            \lccode`\~`\.\lowercase{\def~}{\discretionary{\hbox{\char`\.}}{\Wrappedafterbreak}{\hbox{\char`\.}}}%
            \lccode`\~`\,\lowercase{\def~}{\discretionary{\hbox{\char`\,}}{\Wrappedafterbreak}{\hbox{\char`\,}}}%
            \lccode`\~`\;\lowercase{\def~}{\discretionary{\hbox{\char`\;}}{\Wrappedafterbreak}{\hbox{\char`\;}}}%
            \lccode`\~`\:\lowercase{\def~}{\discretionary{\hbox{\char`\:}}{\Wrappedafterbreak}{\hbox{\char`\:}}}%
            \lccode`\~`\?\lowercase{\def~}{\discretionary{\hbox{\char`\?}}{\Wrappedafterbreak}{\hbox{\char`\?}}}%
            \lccode`\~`\!\lowercase{\def~}{\discretionary{\hbox{\char`\!}}{\Wrappedafterbreak}{\hbox{\char`\!}}}%
            \lccode`\~`\/\lowercase{\def~}{\discretionary{\hbox{\char`\/}}{\Wrappedafterbreak}{\hbox{\char`\/}}}%
            \catcode`\.\active
            \catcode`\,\active
            \catcode`\;\active
            \catcode`\:\active
            \catcode`\?\active
            \catcode`\!\active
            \catcode`\/\active
            \lccode`\~`\~
        }
    \makeatother

    \let\OriginalVerbatim=\Verbatim
    \makeatletter
    \renewcommand{\Verbatim}[1][1]{%
        %\parskip\z@skip
        \sbox\Wrappedcontinuationbox {\Wrappedcontinuationsymbol}%
        \sbox\Wrappedvisiblespacebox {\FV@SetupFont\Wrappedvisiblespace}%
        \def\FancyVerbFormatLine ##1{\hsize\linewidth
            \vtop{\raggedright\hyphenpenalty\z@\exhyphenpenalty\z@
                \doublehyphendemerits\z@\finalhyphendemerits\z@
                \strut ##1\strut}%
        }%
        % If the linebreak is at a space, the latter will be displayed as visible
        % space at end of first line, and a continuation symbol starts next line.
        % Stretch/shrink are however usually zero for typewriter font.
        \def\FV@Space {%
            \nobreak\hskip\z@ plus\fontdimen3\font minus\fontdimen4\font
            \discretionary{\copy\Wrappedvisiblespacebox}{\Wrappedafterbreak}
            {\kern\fontdimen2\font}%
        }%

        % Allow breaks at special characters using \PYG... macros.
        \Wrappedbreaksatspecials
        % Breaks at punctuation characters . , ; ? ! and / need catcode=\active
        \OriginalVerbatim[#1,codes*=\Wrappedbreaksatpunct]%
    }
    \makeatother

    % Exact colors from NB
    \definecolor{incolor}{HTML}{303F9F}
    \definecolor{outcolor}{HTML}{D84315}
    \definecolor{cellborder}{HTML}{CFCFCF}
    \definecolor{cellbackground}{HTML}{F7F7F7}

    % prompt
    \makeatletter
    \newcommand{\boxspacing}{\kern\kvtcb@left@rule\kern\kvtcb@boxsep}
    \makeatother
    \newcommand{\prompt}[4]{
        {\ttfamily\llap{{\color{#2}[#3]:\hspace{3pt}#4}}\vspace{-\baselineskip}}
    }
    

    
    % Prevent overflowing lines due to hard-to-break entities
    \sloppy
    % Setup hyperref package
    \hypersetup{
      breaklinks=true,  % so long urls are correctly broken across lines
      colorlinks=true,
      urlcolor=urlcolor,
      linkcolor=linkcolor,
      citecolor=citecolor,
      }
    % Slightly bigger margins than the latex defaults
    
    \geometry{verbose,tmargin=1in,bmargin=1in,lmargin=1in,rmargin=1in}
    
    

\begin{document}
    
    \maketitle
    
    

    
    \section{Question 1}\label{question-1}

\subsection{Supervised learning in a broad
context}\label{supervised-learning-in-a-broad-context}

\subsubsection{Introduction}\label{introduction}

According to Hastie et al.~(2009), there are three types of statistical
learning problems, namely supervised, unsupervised and semi-supervised
learning. Supervised learning algorithms aim to predict one or more
quantities, referred to as the response value (vector)
\(\mathbf{Y}\in\mathbb{R}^q\), using variables
\(\mathbf{X}\in\mathbb{R}^p\), also called features. An example of such
a learning problem is predicting the probability of a patient having a
heart attack based on features such as their heart rate and blood
pressure. In contrast, unsupervised learning algorithms only observe
features and aim to learn how data is organised or clustered. Clustering
methods can, for example, be applied to clinical data to identify groups
of patients with different survival rates. As the name, semi-supervised,
suggests, this learning problem combines supervised and unsupervised
learning, with features present but only some observations have
corresponding responses. Supervised learning will be the focus for the
remainder of this chapter.

    \subsubsection{The probabilistic approach to supervised
learning}\label{the-probabilistic-approach-to-supervised-learning}

As stated previously, the objective of supervised statistical learning
algorithms is to predict some unknown response \(\mathbf{Y}\) from the
corresponding features \(\mathbf{X}\) with some general real-valued
function \(f:\mathcal{X}\rightarrow \mathbb{R}^p\) referred to as the
candidate model. This model is estimated by fitting it to observed data,
\(D=\{(\mathbf{x}_i,\mathbf{y}_i):i=1,...N\}\) where N denotes the
number of observations, \(\mathbf{x}_i\in\mathcal{X}\in\mathbb{R}^p\)
and \(\mathbf{x}_i\in\mathcal{X}\in\mathbb{R}^p\). Hastie et al.~(2009)
claims this objective is achieved by assuming that
\((\mathbf{X},\mathbf{Y})\) are random vectors with a probability
density function
\(\ddot{\mathbb{P}}_{\mathbf{X,Y}}:\mathcal{X}.\mathcal{Y}\rightarrow[0,\infty)\)
such that the observations in D are identically, independently
distributed according to \(\ddot{\mathbb{P}}\). It follows that
supervised learning methods, therefore, try to fit a model,
\(\hat{f}(\mathbf{x})\), to the observed data, \(D\), to predict
\(\mathbf{Y}\) with \(\hat{f}(\mathbf{X})\).

Depending on the response data type, different supervised learning
problems arise. When \(\mathbf{Y}\) is numeric, predictions are made
with regression and more specifically, if \(q>1\), the problem becomes
multiple output regression. When \(\mathbf{Y}\) is categorical, the aim
is to solve a classification problem. As with regression, classification
can be branched out further into binary, multi-class and multi-label
classification. Binary classification problems predict a binary
categorical variable, \(Y\in\{0,1\}\). When Y needs to be classified
into one of \(k>2\) classes, the problem is multi-class classification,
but when the response can be classified into more than one class and a
vector of categorical variables is therefore predicted, the problem is
multi-label classification. Hastie et al.~(2009)

To illustrate the difference between these supervised learning problems,
consider, for example, the medical domain in which various predictions
can be made. If the probability of a single patient getting infected
with a disease should be predicted, the predicted value will be a
continuous number between zero and one and therefore, this is a
regression problem. When one wants to predict whether a patient is
infected or not, the predicted outcomes can be encoded as zero if the
patient is not infected and one otherwise. This is therefore a binary
classification problem. When a patient is suffering from one of ten
possible diseases, and the aim is to predict which disease the patient
is suffering from, this becomes a multi-class classification problem
since the response is categorical with more than two possible outcomes.

    \subsubsection{The role of loss
functions}\label{the-role-of-loss-functions}

Hastie et al.~(2009) claims that supervised learning problems can be
solved with discriminative or generative modelling. The generative
approach aims to estimate the distribution of \(\mathbf{Y}|\mathbf{X}\)
by first estimating \(\ddot{\mathbb{P}}_{\mathbf{X,Y}}\). Finally, the
estimated distribution of \(\mathbf{Y}|\mathbf{X}\) is then used for
prediction. In contrast, discriminative modelling focuses on estimating
\(\ddot{\mathbb{P}}_{\mathbf{X}}\) nonparametrically in order to extract
information from the distribution \(\mathbf{Y}|\mathbf{X}\) to
ultimately make predictions. The discriminative approach, therefore,
estimates the distribution of \(\mathbf{Y}|\mathbf{X}\) directly by
training a model, \(\hat{f}\), using the structure of \(\mathbf{X}\).

There are various models that can be used to predict \(\mathbf{Y}\).
According to Hastie et al.~(2009), discriminative modelling defines the
optimal model as \(\ddot{f}=argmin_{f}[\ddot{\mathcal{R}}(f)]\).
Therefore, the optimal model, also known in this setting as the
population minimiser, minimises the population risk,
\(\ddot{\mathcal{R}}(f)=E_{\mathbf{X, Y}}[\mathcal{L}(\mathbf{Y},f(\mathbf{X}))]\).
In this formulation \(\mathcal{L}(\mathbf{Y},f(\mathbf{X}))\) is some
loss function
\(\mathcal{L}:\mathcal{Y}\times\mathbb{R}^r\rightarrow[0,\infty)\) that
quantifies the loss occured when \(f(\mathbf{X})\) is used to predict
\(\mathbf{Y}\). Equivalently, the model that minimises the loss function
will be the best model.

Just as different models can be found to fit data, different loss
functions can be employed to find the best model. Intuitively, different
loss functions will yield different optimal models. To show this,
consider two well-known loss functions, the squared-error loss,
\(\mathcal{L}(Y,f(\mathbf{X}))=(Y-f(\mathbf{X}))^2\), and absolute loss,
\(\mathcal{L}(Y,f(\mathbf{X}))=|Y-f(\mathbf{X})|\). The optimal model
for \(\mathbf{X}=\mathbf{x}\) using the squared-error loss is derived as
Hastie et al.~(2009):
\begin{align*}
  \ddot{f}(\mathbf{x})&=argmin_{f}[E((Y-f(\mathbf{X}))^2|\mathbf{X}=\mathbf{x})] \\
  &=argmin_{f}[E{(Y^2|\mathbf{X}=\mathbf{x})-2f(\mathbf{x})E(Y|\mathbf{X}=\mathbf{x})+f(\mathbf{x})^2}] \\
\therefore \text{ The expression }&\text{will be minimized when } f(\mathbf{x})E(Y|\mathbf{X}=\mathbf{x})=f(\mathbf{x})^2 \\
\therefore& \ddot{f}(\mathbf{x})= E(Y|\mathbf{X}=\mathbf{x}) \\
\end{align*}

It follows that a similar procedure can be followed for the absolute
loss, and although the calculations are not as simple, it can be derived
that the population minimiser at \(\mathbf{X=x}\) is equal to the median
of \(\mathbf{Y}|\mathbf{X=x}\). It is therefore shown that two different
loss functions result in different optimal models. The squared error
loss is simpler to work with, but it is less robust than the absolute
loss. Loss functions should therefore be carefully chosen after
considering the nature of the supervised learning problem.

Another common loss function, specifically used in classification, is
the 0-1 loss, \[
\mathcal{L}_{0-1}(Y,f(\mathbf{X}))=\begin{cases}
                                0 &\text{if } Y=f(\mathbf{X}), \\
                                1 &\text{if } Y\ne f(\mathbf{X}). 
                              \end{cases}
\]

The derivation for finding the corresponding optimal model follows as:
\begin{align*}
  \ddot{f}(\mathbf{x})&=argmin_{f}[E(\mathbb{I}(Y\ne f(\mathbf{x}))|\mathbf{X}=\mathbf{x})] \\
  &= argmin_{f}[p(\mathbf{x})\mathbb{I}(1 \ne f(\mathbf{x}))+(1-p(\mathbf{x}))\mathbb{I}(0 \ne f(\mathbf{x}))] \\
  &= argmin_{f}[p(\mathbf{x})\mathbb{I}(0 = f(\mathbf{x}))+(1-p(\mathbf{x}))(1-\mathbb{I}(0 = f(\mathbf{x})))] \\
  &= argmin_{f}\begin{cases}
                    p(\mathbf{x}) &\text{if } f(\mathbf{x})=0, \\
                    1-p(\mathbf{x}) &\text{if } f(\mathbf{x})=1. \\
               \end{cases}\\
\therefore \ddot{f}(\mathbf{x})&=\begin{cases}
                                0 &\text{if } p(\mathbf{x})<0.5, \\
                                1 &\text{if } p(\mathbf{x})\ge 0.5. 
                              \end{cases}
\end{align*}

Since this loss function is used in classification problems, the optimal
model is not a spesific quantity as with the squared error loss and
absolute error loss, but rather a decision boundary. To show how this
would look visualy, this 0-1 loss function was applied to a logistic
regression problem which required patients to be classified as heart
attack risk patients (Y=1) or not (Y=0) based on their age (\(X_1\)) and
cholestrol levels (\(X_2\)). The first step in this application is to
simulate a sample of size n = 200 from the true data generating
mechanism, \(Y\sim Ber(p)\) with \(p=[1+e^{\mathbf{-X\beta}}]^{-1}\)
where \(\mathbf{\beta}'=[\beta_0,\beta_1,\beta_2]=[-2,0.02,0.01]\),
\(\mathbf{X}'=[X_1,X_2]\), \(X_1\sim N(50,10)\) and
\(X_2\sim N(200,30)\).

    \begin{tcolorbox}[breakable, size=fbox, boxrule=1pt, pad at break*=1mm,colback=cellbackground, colframe=cellborder]
\prompt{In}{incolor}{1}{\boxspacing}
\begin{Verbatim}[commandchars=\\\{\}]
\PY{k+kn}{import}\PY{+w}{ }\PY{n+nn}{warnings}
\PY{n}{warnings}\PY{o}{.}\PY{n}{filterwarnings}\PY{p}{(}\PY{l+s+s1}{\PYZsq{}}\PY{l+s+s1}{ignore}\PY{l+s+s1}{\PYZsq{}}\PY{p}{)}
\end{Verbatim}
\end{tcolorbox}

    \begin{tcolorbox}[breakable, size=fbox, boxrule=1pt, pad at break*=1mm,colback=cellbackground, colframe=cellborder]
\prompt{In}{incolor}{2}{\boxspacing}
\begin{Verbatim}[commandchars=\\\{\}]
\PY{k+kn}{import}\PY{+w}{ }\PY{n+nn}{numpy}\PY{+w}{ }\PY{k}{as}\PY{+w}{ }\PY{n+nn}{np}
\PY{k+kn}{import}\PY{+w}{ }\PY{n+nn}{matplotlib}\PY{n+nn}{.}\PY{n+nn}{pyplot}\PY{+w}{ }\PY{k}{as}\PY{+w}{ }\PY{n+nn}{plt}
      
\PY{n}{Beta} \PY{o}{=} \PY{p}{[}\PY{o}{\PYZhy{}}\PY{l+m+mi}{2}\PY{p}{,} \PY{l+m+mf}{0.02}\PY{p}{,} \PY{l+m+mf}{0.01}\PY{p}{]}      

\PY{k}{def}\PY{+w}{ }\PY{n+nf}{sigmoid}\PY{p}{(}\PY{n}{z}\PY{p}{)}\PY{p}{:}
    \PY{k}{return} \PY{l+m+mi}{1} \PY{o}{/} \PY{p}{(}\PY{l+m+mi}{1} \PY{o}{+} \PY{n}{np}\PY{o}{.}\PY{n}{exp}\PY{p}{(}\PY{o}{\PYZhy{}}\PY{n}{z}\PY{p}{)}\PY{p}{)}

\PY{k}{def}\PY{+w}{ }\PY{n+nf}{generate\PYZus{}data}\PY{p}{(}\PY{n}{n}\PY{p}{)}\PY{p}{:}
    \PY{n}{age}\PY{o}{=} \PY{n}{np}\PY{o}{.}\PY{n}{random}\PY{o}{.}\PY{n}{normal}\PY{p}{(}\PY{n}{loc}\PY{o}{=}\PY{l+m+mi}{50}\PY{p}{,} \PY{n}{scale}\PY{o}{=}\PY{l+m+mi}{10}\PY{p}{,} \PY{n}{size}\PY{o}{=}\PY{n}{n}\PY{p}{)}           
    \PY{n}{cholesterol} \PY{o}{=} \PY{n}{np}\PY{o}{.}\PY{n}{random}\PY{o}{.}\PY{n}{normal}\PY{p}{(}\PY{n}{loc}\PY{o}{=}\PY{l+m+mi}{200}\PY{p}{,} \PY{n}{scale}\PY{o}{=}\PY{l+m+mi}{30}\PY{p}{,} \PY{n}{size}\PY{o}{=}\PY{n}{n}\PY{p}{)}  
    \PY{n}{X} \PY{o}{=} \PY{n}{np}\PY{o}{.}\PY{n}{column\PYZus{}stack}\PY{p}{(}\PY{p}{(}\PY{n}{age}\PY{p}{,} \PY{n}{cholesterol}\PY{p}{)}\PY{p}{)}
    \PY{n}{X\PYZus{}design} \PY{o}{=} \PY{n}{np}\PY{o}{.}\PY{n}{c\PYZus{}}\PY{p}{[}\PY{n}{np}\PY{o}{.}\PY{n}{ones}\PY{p}{(}\PY{n}{X}\PY{o}{.}\PY{n}{shape}\PY{p}{[}\PY{l+m+mi}{0}\PY{p}{]}\PY{p}{)}\PY{p}{,} \PY{n}{X}\PY{p}{]} 
    \PY{n}{z} \PY{o}{=} \PY{n}{X\PYZus{}design} \PY{o}{@} \PY{n}{Beta}
    \PY{n}{p} \PY{o}{=} \PY{n}{sigmoid}\PY{p}{(}\PY{n}{z}\PY{p}{)}
    \PY{n}{y} \PY{o}{=} \PY{n}{np}\PY{o}{.}\PY{n}{random}\PY{o}{.}\PY{n}{binomial}\PY{p}{(}\PY{l+m+mi}{1}\PY{p}{,} \PY{n}{p}\PY{p}{)}

    \PY{k}{return} \PY{n}{X}\PY{p}{,} \PY{n}{y}

\PY{n}{np}\PY{o}{.}\PY{n}{random}\PY{o}{.}\PY{n}{seed}\PY{p}{(}\PY{l+m+mi}{123}\PY{p}{)}
\PY{n}{X}\PY{p}{,} \PY{n}{y} \PY{o}{=} \PY{n}{generate\PYZus{}data}\PY{p}{(}\PY{l+m+mi}{1000}\PY{p}{)}
\end{Verbatim}
\end{tcolorbox}

    Next a heatmap was contructed to visualise
\(E(\mathcal{L}_{0-1}(Y,f(\mathbf{X})))\) for varying values of
\(\beta_1\) and \(\beta_2\) while keeping \(\beta_0\) fixed.

    \begin{tcolorbox}[breakable, size=fbox, boxrule=1pt, pad at break*=1mm,colback=cellbackground, colframe=cellborder]
\prompt{In}{incolor}{3}{\boxspacing}
\begin{Verbatim}[commandchars=\\\{\}]
\PY{k+kn}{import}\PY{+w}{ }\PY{n+nn}{numpy}\PY{+w}{ }\PY{k}{as}\PY{+w}{ }\PY{n+nn}{np}
\PY{k+kn}{import}\PY{+w}{ }\PY{n+nn}{matplotlib}\PY{n+nn}{.}\PY{n+nn}{pyplot}\PY{+w}{ }\PY{k}{as}\PY{+w}{ }\PY{n+nn}{plt}

\PY{n}{b1} \PY{o}{=} \PY{n}{np}\PY{o}{.}\PY{n}{linspace}\PY{p}{(}\PY{o}{\PYZhy{}}\PY{l+m+mi}{1}\PY{p}{,} \PY{l+m+mi}{1}\PY{p}{,} \PY{l+m+mi}{500}\PY{p}{)}
\PY{n}{b2} \PY{o}{=} \PY{n}{np}\PY{o}{.}\PY{n}{linspace}\PY{p}{(}\PY{o}{\PYZhy{}}\PY{l+m+mi}{1}\PY{p}{,} \PY{l+m+mi}{1}\PY{p}{,} \PY{l+m+mi}{500}\PY{p}{)}
\PY{n}{B1}\PY{p}{,} \PY{n}{B2} \PY{o}{=} \PY{n}{np}\PY{o}{.}\PY{n}{meshgrid}\PY{p}{(}\PY{n}{b1}\PY{p}{,} \PY{n}{b2}\PY{p}{)}
\PY{n}{X\PYZus{}design} \PY{o}{=} \PY{n}{np}\PY{o}{.}\PY{n}{c\PYZus{}}\PY{p}{[}\PY{n}{np}\PY{o}{.}\PY{n}{ones}\PY{p}{(}\PY{n}{X}\PY{o}{.}\PY{n}{shape}\PY{p}{[}\PY{l+m+mi}{0}\PY{p}{]}\PY{p}{)}\PY{p}{,} \PY{n}{X}\PY{p}{]} 

\PY{n}{loss} \PY{o}{=} \PY{n}{np}\PY{o}{.}\PY{n}{zeros\PYZus{}like}\PY{p}{(}\PY{n}{B1}\PY{p}{)}
\PY{k}{for} \PY{n}{i} \PY{o+ow}{in} \PY{n+nb}{range}\PY{p}{(}\PY{n}{B1}\PY{o}{.}\PY{n}{shape}\PY{p}{[}\PY{l+m+mi}{0}\PY{p}{]}\PY{p}{)}\PY{p}{:}
    \PY{k}{for} \PY{n}{j} \PY{o+ow}{in} \PY{n+nb}{range}\PY{p}{(}\PY{n}{B1}\PY{o}{.}\PY{n}{shape}\PY{p}{[}\PY{l+m+mi}{1}\PY{p}{]}\PY{p}{)}\PY{p}{:}
        \PY{n}{beta} \PY{o}{=} \PY{n}{np}\PY{o}{.}\PY{n}{array}\PY{p}{(}\PY{p}{[}\PY{o}{\PYZhy{}}\PY{l+m+mi}{2}\PY{p}{,} \PY{n}{B1}\PY{p}{[}\PY{n}{i}\PY{p}{,} \PY{n}{j}\PY{p}{]}\PY{p}{,} \PY{n}{B2}\PY{p}{[}\PY{n}{i}\PY{p}{,} \PY{n}{j}\PY{p}{]}\PY{p}{]}\PY{p}{)}
        \PY{n}{p\PYZus{}model} \PY{o}{=} \PY{n}{sigmoid}\PY{p}{(}\PY{n}{X\PYZus{}design} \PY{o}{@} \PY{n}{beta}\PY{p}{)}
        \PY{n}{yhat} \PY{o}{=} \PY{p}{(}\PY{n}{p\PYZus{}model} \PY{o}{\PYZgt{}}\PY{o}{=} \PY{l+m+mf}{0.5}\PY{p}{)}\PY{o}{.}\PY{n}{astype}\PY{p}{(}\PY{n+nb}{int}\PY{p}{)}
        \PY{n}{loss}\PY{p}{[}\PY{n}{i}\PY{p}{,} \PY{n}{j}\PY{p}{]} \PY{o}{=} \PY{n}{np}\PY{o}{.}\PY{n}{mean}\PY{p}{(}\PY{n}{yhat} \PY{o}{!=} \PY{n}{y}\PY{p}{)}
\end{Verbatim}
\end{tcolorbox}

    \begin{tcolorbox}[breakable, size=fbox, boxrule=1pt, pad at break*=1mm,colback=cellbackground, colframe=cellborder]
\prompt{In}{incolor}{4}{\boxspacing}
\begin{Verbatim}[commandchars=\\\{\}]
\PY{n}{plt}\PY{o}{.}\PY{n}{figure}\PY{p}{(}\PY{n}{figsize}\PY{o}{=}\PY{p}{(}\PY{l+m+mi}{8}\PY{p}{,}\PY{l+m+mi}{6}\PY{p}{)}\PY{p}{)}
\PY{n}{plt}\PY{o}{.}\PY{n}{imshow}\PY{p}{(}\PY{n}{loss}\PY{p}{,} \PY{n}{origin}\PY{o}{=}\PY{l+s+s1}{\PYZsq{}}\PY{l+s+s1}{lower}\PY{l+s+s1}{\PYZsq{}}\PY{p}{,} \PY{n}{extent}\PY{o}{=}\PY{p}{[}\PY{n}{b1}\PY{o}{.}\PY{n}{min}\PY{p}{(}\PY{p}{)}\PY{p}{,} \PY{n}{b1}\PY{o}{.}\PY{n}{max}\PY{p}{(}\PY{p}{)}\PY{p}{,} \PY{n}{b2}\PY{o}{.}\PY{n}{min}\PY{p}{(}\PY{p}{)}\PY{p}{,} \PY{n}{b2}\PY{o}{.}\PY{n}{max}\PY{p}{(}\PY{p}{)}\PY{p}{]}\PY{p}{,} \PY{n}{aspect}\PY{o}{=}\PY{l+s+s1}{\PYZsq{}}\PY{l+s+s1}{auto}\PY{l+s+s1}{\PYZsq{}}\PY{p}{)}
\PY{n}{plt}\PY{o}{.}\PY{n}{plot}\PY{p}{(}\PY{l+m+mf}{0.02}\PY{p}{,}\PY{l+m+mf}{0.01}\PY{p}{,}\PY{l+s+s2}{\PYZdq{}}\PY{l+s+s2}{ro}\PY{l+s+s2}{\PYZdq{}}\PY{p}{)}
\PY{n}{plt}\PY{o}{.}\PY{n}{colorbar}\PY{p}{(}\PY{n}{label}\PY{o}{=}\PY{l+s+s2}{\PYZdq{}}\PY{l+s+s2}{Expected 0\PYZhy{}1 Loss}\PY{l+s+s2}{\PYZdq{}}\PY{p}{)}
\PY{n}{plt}\PY{o}{.}\PY{n}{xlabel}\PY{p}{(}\PY{l+s+s2}{\PYZdq{}}\PY{l+s+s2}{B1}\PY{l+s+s2}{\PYZdq{}}\PY{p}{)}
\PY{n}{plt}\PY{o}{.}\PY{n}{ylabel}\PY{p}{(}\PY{l+s+s2}{\PYZdq{}}\PY{l+s+s2}{B2}\PY{l+s+s2}{\PYZdq{}}\PY{p}{)}
\PY{n}{plt}\PY{o}{.}\PY{n}{title}\PY{p}{(}\PY{l+s+s2}{\PYZdq{}}\PY{l+s+s2}{Expected 0\PYZhy{}1 Loss Surface}\PY{l+s+s2}{\PYZdq{}}\PY{p}{)}
\PY{n}{plt}\PY{o}{.}\PY{n}{show}\PY{p}{(}\PY{p}{)}
\end{Verbatim}
\end{tcolorbox}

    \begin{center}
    \adjustimage{max size={0.9\linewidth}{0.9\paperheight}}{Ass1_1_files/Ass1_1_7_0.png}
    \end{center}
    { \hspace*{\fill} \\}
    
    As expcted, the heatmap does not show
\(E(\mathcal{L}_{0-1}(Y,f(\mathbf{X})))\) to have a minimum at a certain
point but rather indicates a decision boundary at the true parameter
values. This aligns with what was derived analytically.

    \subsubsection{Empirical risk and
regularisation}\label{empirical-risk-and-regularisation}

In the previous section, it was discussed how to find the optimal model
for a supervised learning problem. When we only consider models in the
model space \(\mathcal{F}\), \(\ddot{f}\) can be approximated with
\(\dot{f}=argmin_{f\in \mathcal{F}}[\ddot{\mathcal{R}}(f)]\). To find a
\(\dot{f}\), however, the data from the whole population or the true
data-generating mechanism is required. Since this is almost never
available in practice, the population minimiser can be approximated with
an optimal model,\(\hat{f}\) that minimises the empirical risk,
\(\hat{\mathcal{R}}(f,\mathcal{D})=\frac{1}{|D|}\sum_{(\mathbf{(x,y)})\in \mathcal{D}}\mathcal{L}(\mathbf{y},f(\mathbf{x}))\).
This is equivalent to the optimasation problem
\[\hat{f}(.,\mathcal{D})=argmin_{f\in \mathcal{F}}[\hat{\mathcal{R}}(f,\mathcal{D})] \tag{1}\]

The empirical risk is a sufficient estimator for the population risk
since
\(\hat{\mathcal{R}}(f,\mathcal{D})\overset{p}{\rightarrow}\ddot{\mathcal{R}}(f)\)
as \(|D|\rightarrow\infty\). To ensure that \(\hat{f}(.,\mathcal{D})\)
is also a sufficient estimator for \(\dot{f}\) the following must hold:
\[\text{sup}_{f\in \mathcal{F}}[|\hat{\mathcal{R}}(f,\mathcal{D})-\ddot{\mathcal{R}}(f)|]\overset{p}{\rightarrow}0\text{ as }|\mathcal{D}|\rightarrow\infty\]
This condition can only be achieved when suitable constraints are
imposed on \(\mathcal{F}\), such as regularisation. According to Hastie
et al.~(2009), regularisation is a learning algorithm that partially
minimises the empirical risk. Regularisation introduces bias into the
model by reducing model complexity. This will enhance model performance
on unseen data, or, differently stated, it will improve model
generalisation.

Types of regularization techniques include variable selection and
functional regularization. Variable selection reduces model complexity
by excluding certain features from the feature set that is used to train
the model, to potentially remove noise and improve model generalization.
This approach is referred to as hard regularization. In contrast,
functional regularization reduces model complexity by introducing a
penalty function \(\mathcal{P:F}\rightarrow [0,\infty)\) which changes
the model space to
\(\mathcal{F}(t)=\{f\in \mathcal{F}:\mathcal{P}\le t\}\). With
functional regularization, optimisation problem (1) therefore becomes:
\[\hat{f}(.,\mathcal{D},\lambda)=argmin_{f\in \mathcal{F}}[\hat{\mathcal{R}}(f)+\lambda \mathcal{P}(f)] \]

Continuing the logistic regression example from before, functional
regularisation, spesifically ridge regularisation, can be applied to
potentialy improve generalisation performance. For ridge regularisation
the penalty function becomes
\(\mathcal{P}(f)=||E_{\mathbf{X}}[\frac{\partial f}{\partial \mathbf{X}}]||^2_2\)
and when applied to logistic regression this reduces to
\(\mathbf{\beta}'\mathbf{\beta}\). The following implementation
performed logistic regression with and without regularisation while
setting the tuning parameter \(\lambda=1\).

    \begin{tcolorbox}[breakable, size=fbox, boxrule=1pt, pad at break*=1mm,colback=cellbackground, colframe=cellborder]
\prompt{In}{incolor}{5}{\boxspacing}
\begin{Verbatim}[commandchars=\\\{\}]
\PY{k+kn}{from}\PY{+w}{ }\PY{n+nn}{sklearn}\PY{n+nn}{.}\PY{n+nn}{linear\PYZus{}model}\PY{+w}{ }\PY{k+kn}{import} \PY{n}{LogisticRegression}
\PY{k+kn}{from}\PY{+w}{ }\PY{n+nn}{sklearn}\PY{n+nn}{.}\PY{n+nn}{preprocessing}\PY{+w}{ }\PY{k+kn}{import} \PY{n}{StandardScaler}

\PY{k}{def}\PY{+w}{ }\PY{n+nf}{train\PYZus{}models}\PY{p}{(}\PY{p}{)}\PY{p}{:}
    \PY{n}{model\PYZus{}unreg} \PY{o}{=} \PY{n}{LogisticRegression}\PY{p}{(} \PY{n}{C}\PY{o}{=}\PY{n}{np}\PY{o}{.}\PY{n}{inf}\PY{p}{,}\PY{n}{fit\PYZus{}intercept}\PY{o}{=}\PY{k+kc}{True}\PY{p}{,} \PY{n}{random\PYZus{}state}\PY{o}{=}\PY{l+m+mi}{0}\PY{p}{,} \PY{n}{max\PYZus{}iter}\PY{o}{=}\PY{l+m+mi}{1000}\PY{p}{)}
    \PY{n}{model\PYZus{}wridge} \PY{o}{=} \PY{n}{LogisticRegression}\PY{p}{(} \PY{n}{C}\PY{o}{=}\PY{l+m+mi}{1}\PY{p}{,} \PY{n}{fit\PYZus{}intercept}\PY{o}{=}\PY{k+kc}{True}\PY{p}{,} \PY{n}{random\PYZus{}state}\PY{o}{=}\PY{l+m+mi}{0}\PY{p}{,} \PY{n}{max\PYZus{}iter}\PY{o}{=}\PY{l+m+mi}{1000}\PY{p}{)}
    \PY{n}{model\PYZus{}sridge} \PY{o}{=} \PY{n}{LogisticRegression}\PY{p}{(} \PY{n}{C}\PY{o}{=}\PY{l+m+mf}{0.005}\PY{p}{,} \PY{n}{fit\PYZus{}intercept}\PY{o}{=}\PY{k+kc}{True}\PY{p}{,} \PY{n}{random\PYZus{}state}\PY{o}{=}\PY{l+m+mi}{0}\PY{p}{,} \PY{n}{max\PYZus{}iter}\PY{o}{=}\PY{l+m+mi}{1000}\PY{p}{)}

    \PY{k}{return} \PY{n}{model\PYZus{}unreg}\PY{p}{,} \PY{n}{model\PYZus{}wridge}\PY{p}{,} \PY{n}{model\PYZus{}sridge}

\PY{n}{scaler} \PY{o}{=} \PY{n}{StandardScaler}\PY{p}{(}\PY{p}{)}
\PY{n}{X\PYZus{}scaled} \PY{o}{=} \PY{n}{scaler}\PY{o}{.}\PY{n}{fit\PYZus{}transform}\PY{p}{(}\PY{n}{X}\PY{p}{)}

\PY{n}{model\PYZus{}unreg}\PY{p}{,} \PY{n}{model\PYZus{}wridge}\PY{p}{,} \PY{n}{model\PYZus{}sridge} \PY{o}{=} \PY{n}{train\PYZus{}models}\PY{p}{(}\PY{p}{)}
\PY{n}{model\PYZus{}unreg}\PY{o}{.}\PY{n}{fit}\PY{p}{(}\PY{n}{X\PYZus{}scaled}\PY{p}{,} \PY{n}{y}\PY{p}{)}
\PY{n}{model\PYZus{}wridge}\PY{o}{.}\PY{n}{fit}\PY{p}{(}\PY{n}{X\PYZus{}scaled}\PY{p}{,} \PY{n}{y}\PY{p}{)}
\PY{n}{model\PYZus{}sridge}\PY{o}{.}\PY{n}{fit}\PY{p}{(}\PY{n}{X\PYZus{}scaled}\PY{p}{,} \PY{n}{y}\PY{p}{)}

\PY{n+nb}{print}\PY{p}{(}\PY{l+s+s2}{\PYZdq{}}\PY{l+s+s2}{Unregularised Logistic Regression}\PY{l+s+s2}{\PYZdq{}}\PY{p}{)}
\PY{n+nb}{print}\PY{p}{(}\PY{l+s+s2}{\PYZdq{}}\PY{l+s+s2}{Intercept   : }\PY{l+s+s2}{\PYZdq{}}\PY{p}{,}\PY{n}{model\PYZus{}unreg}\PY{o}{.}\PY{n}{intercept\PYZus{}}\PY{p}{)}
\PY{n+nb}{print}\PY{p}{(}\PY{l+s+s2}{\PYZdq{}}\PY{l+s+s2}{Coefficients: }\PY{l+s+s2}{\PYZdq{}}\PY{p}{,}\PY{n}{model\PYZus{}unreg}\PY{o}{.}\PY{n}{coef\PYZus{}}\PY{p}{[}\PY{l+m+mi}{0}\PY{p}{]}\PY{p}{)}

\PY{n+nb}{print}\PY{p}{(}\PY{l+s+s2}{\PYZdq{}}\PY{l+s+se}{\PYZbs{}n}\PY{l+s+s2}{Ridge Logistic Regression with C = 1}\PY{l+s+s2}{\PYZdq{}}\PY{p}{)}
\PY{n+nb}{print}\PY{p}{(}\PY{l+s+s2}{\PYZdq{}}\PY{l+s+s2}{Intercept   : }\PY{l+s+s2}{\PYZdq{}}\PY{p}{,}\PY{n}{model\PYZus{}wridge}\PY{o}{.}\PY{n}{intercept\PYZus{}}\PY{p}{)}
\PY{n+nb}{print}\PY{p}{(}\PY{l+s+s2}{\PYZdq{}}\PY{l+s+s2}{Coefficients: }\PY{l+s+s2}{\PYZdq{}}\PY{p}{,}\PY{n}{model\PYZus{}wridge}\PY{o}{.}\PY{n}{coef\PYZus{}}\PY{p}{[}\PY{l+m+mi}{0}\PY{p}{]}\PY{p}{)}

\PY{n+nb}{print}\PY{p}{(}\PY{l+s+s2}{\PYZdq{}}\PY{l+s+se}{\PYZbs{}n}\PY{l+s+s2}{Ridge Logistic Regression with C = 0.005}\PY{l+s+s2}{\PYZdq{}}\PY{p}{)}
\PY{n+nb}{print}\PY{p}{(}\PY{l+s+s2}{\PYZdq{}}\PY{l+s+s2}{Intercept   : }\PY{l+s+s2}{\PYZdq{}}\PY{p}{,}\PY{n}{model\PYZus{}sridge}\PY{o}{.}\PY{n}{intercept\PYZus{}}\PY{p}{)}
\PY{n+nb}{print}\PY{p}{(}\PY{l+s+s2}{\PYZdq{}}\PY{l+s+s2}{Coefficients: }\PY{l+s+s2}{\PYZdq{}}\PY{p}{,}\PY{n}{model\PYZus{}sridge}\PY{o}{.}\PY{n}{coef\PYZus{}}\PY{p}{[}\PY{l+m+mi}{0}\PY{p}{]}\PY{p}{)}
\end{Verbatim}
\end{tcolorbox}

    \begin{Verbatim}[commandchars=\\\{\}]
Unregularised Logistic Regression
Intercept   :  [0.97870688]
Coefficients:  [0.326      0.33730111]

Ridge Logistic Regression with C = 1
Intercept   :  [0.97822383]
Coefficients:  [0.32414595 0.33538198]

Ridge Logistic Regression with C = 0.005
Intercept   :  [0.94557328]
Coefficients:  [0.15615171 0.16186327]
    \end{Verbatim}

    The tuning parameter C plays a significant role in the amount by wich a
model is regularized. When is large, model complexity is penalised less
than when C is chosen small. From the results it is clear that as
regularisation increased, estimates moved closer to the true parameter
values. It is naive however to make assumption about the generalisation
performance of the different models by just comparing the parameter
estimates, which motivates the reason for computing generalisation
errors.

    \subsubsection{Importance of measuring the generalisation
performance}\label{importance-of-measuring-the-generalisation-performance}

When a model is too complex and fits the observed data too well,
overfitting occurs. It captured random fluctuations along with the true
underlying pattern in the data, which causes poor generalisation. In
contrast, underfitting occurs when the model is so simple that it cannot
learn the underlying pattern at all, which will also lead to poor
generalisation. It is therefore important that the generalisation
performance is measured to evaluate a fitted model and to ensure that
neither overfitting nor underfitting occurs.

There are different measurements for generalisation performance. The
generalisation error,
\(Err(\hat{f},\mathcal{D})=E_{\mathbf{X,Y}}[\mathcal{L}(\mathbf{Y},\hat{f}(\mathbf{X},\mathcal{D}))]\),
measures the error of a specific fitted model on unseen data. Since the
true data-generating mechanism is unknown, the error must be estimated.
The training error,
\(\bar{err}(\hat{f},\mathcal{D})=\hat{\mathcal{R}}(\hat{f}(.,\mathcal{D}),\mathcal{D})\)
is considered a downword biased estimator for the generalisation error
since it fits and assess the model with the same data.

A better estimator for the generalisation error is the test error,
\(\hat{Err}(\hat{f},\mathcal{D},\mathcal{T}_e)=\hat{\mathcal{R}}(\hat{f}(.,\mathcal{D}),\mathcal{T}_e)\).
This estimate is obtained by first dividing the available data into a
set \(\mathcal{D}\), for model estimation and a set \(\mathcal{T}_e\),
for performance assessment. One of the test error's attractive
properties is that it becomes consistent as the \(|\mathcal{T}_e|\)
grows if the data-generating mechanism satisfies the i.i.d. assumption.

In contrast the expected generalisation error,
\(Err(\hat{f})=E_{\mathcal{D}}[Err(\hat{f},\mathcal{D})]\) measures the
error of a learning algorithm. This error can be estimated with
simulation. First, generate a large number, say B, data sets from the
data-generating mechanism and split each set into
\((\mathcal{D}^{(b)},\mathcal{T}_e^{(b)})\) for b=1,\ldots,B. It follows
that the expected generalisation error can be estimated with
\(\hat{Err}(\hat{f},B)=\frac{1}{B}\sum^B_{b=1}\hat{Err}(\hat{f},\mathcal{D}^{(b)},\mathcal{T}_e^{(b)})\)

To explore how regularisation influence the generalisation performance
of a model, the training, test and estimated expected generalisation
error can now be computed for the logistic models previously trained.

    \begin{tcolorbox}[breakable, size=fbox, boxrule=1pt, pad at break*=1mm,colback=cellbackground, colframe=cellborder]
\prompt{In}{incolor}{6}{\boxspacing}
\begin{Verbatim}[commandchars=\\\{\}]
\PY{k+kn}{from}\PY{+w}{ }\PY{n+nn}{sklearn}\PY{n+nn}{.}\PY{n+nn}{model\PYZus{}selection}\PY{+w}{ }\PY{k+kn}{import} \PY{n}{train\PYZus{}test\PYZus{}split}
\PY{k+kn}{from}\PY{+w}{ }\PY{n+nn}{sklearn}\PY{n+nn}{.}\PY{n+nn}{metrics}\PY{+w}{ }\PY{k+kn}{import} \PY{n}{zero\PYZus{}one\PYZus{}loss}

\PY{k}{def}\PY{+w}{ }\PY{n+nf}{fit\PYZus{}and\PYZus{}evaluate}\PY{p}{(}\PY{n}{model}\PY{p}{,} \PY{n}{X\PYZus{}train}\PY{p}{,} \PY{n}{y\PYZus{}train}\PY{p}{,} \PY{n}{X\PYZus{}test}\PY{p}{,} \PY{n}{y\PYZus{}test}\PY{p}{)}\PY{p}{:}
    \PY{n}{model}\PY{o}{.}\PY{n}{fit}\PY{p}{(}\PY{n}{X\PYZus{}train}\PY{p}{,} \PY{n}{y\PYZus{}train}\PY{p}{)}
    \PY{n}{train\PYZus{}err} \PY{o}{=} \PY{n}{zero\PYZus{}one\PYZus{}loss}\PY{p}{(}\PY{n}{y\PYZus{}train}\PY{p}{,} \PY{n}{model}\PY{o}{.}\PY{n}{predict}\PY{p}{(}\PY{n}{X\PYZus{}train}\PY{p}{)}\PY{p}{)}
    \PY{n}{test\PYZus{}err} \PY{o}{=} \PY{n}{zero\PYZus{}one\PYZus{}loss}\PY{p}{(}\PY{n}{y\PYZus{}test}\PY{p}{,} \PY{n}{model}\PY{o}{.}\PY{n}{predict}\PY{p}{(}\PY{n}{X\PYZus{}test}\PY{p}{)}\PY{p}{)}
    
    \PY{k}{return} \PY{n}{train\PYZus{}err}\PY{p}{,} \PY{n}{test\PYZus{}err}

\PY{n}{X\PYZus{}train}\PY{p}{,} \PY{n}{X\PYZus{}test}\PY{p}{,} \PY{n}{y\PYZus{}train}\PY{p}{,} \PY{n}{y\PYZus{}test} \PY{o}{=} \PY{n}{train\PYZus{}test\PYZus{}split}\PY{p}{(}\PY{n}{X\PYZus{}scaled}\PY{p}{,} \PY{n}{y}\PY{p}{,} \PY{n}{test\PYZus{}size}\PY{o}{=}\PY{l+m+mf}{0.3}\PY{p}{,} \PY{n}{random\PYZus{}state}\PY{o}{=}\PY{l+m+mi}{0}\PY{p}{)}

\PY{c+c1}{\PYZsh{} Obtain test and training errors}
\PY{n}{model\PYZus{}unreg}\PY{p}{,} \PY{n}{model\PYZus{}wridge}\PY{p}{,} \PY{n}{model\PYZus{}sridge} \PY{o}{=} \PY{n}{train\PYZus{}models}\PY{p}{(}\PY{p}{)}
\PY{n}{train\PYZus{}err\PYZus{}unreg}\PY{p}{,} \PY{n}{test\PYZus{}err\PYZus{}unreg} \PY{o}{=} \PY{n}{fit\PYZus{}and\PYZus{}evaluate}\PY{p}{(}\PY{n}{model\PYZus{}unreg}\PY{p}{,} \PY{n}{X\PYZus{}train}\PY{p}{,} \PY{n}{y\PYZus{}train}\PY{p}{,} \PY{n}{X\PYZus{}test}\PY{p}{,} \PY{n}{y\PYZus{}test}\PY{p}{)}
\PY{n}{train\PYZus{}err\PYZus{}wridge}\PY{p}{,} \PY{n}{test\PYZus{}err\PYZus{}wridge} \PY{o}{=} \PY{n}{fit\PYZus{}and\PYZus{}evaluate}\PY{p}{(}\PY{n}{model\PYZus{}wridge}\PY{p}{,} \PY{n}{X\PYZus{}train}\PY{p}{,} \PY{n}{y\PYZus{}train}\PY{p}{,} \PY{n}{X\PYZus{}test}\PY{p}{,} \PY{n}{y\PYZus{}test}\PY{p}{)}
\PY{n}{train\PYZus{}err\PYZus{}sridge}\PY{p}{,} \PY{n}{test\PYZus{}err\PYZus{}sridge} \PY{o}{=} \PY{n}{fit\PYZus{}and\PYZus{}evaluate}\PY{p}{(}\PY{n}{model\PYZus{}sridge}\PY{p}{,} \PY{n}{X\PYZus{}train}\PY{p}{,} \PY{n}{y\PYZus{}train}\PY{p}{,} \PY{n}{X\PYZus{}test}\PY{p}{,} \PY{n}{y\PYZus{}test}\PY{p}{)}


\PY{c+c1}{\PYZsh{} Obtain expected test errors}
\PY{k}{def}\PY{+w}{ }\PY{n+nf}{estimate\PYZus{}expected\PYZus{}error}\PY{p}{(}\PY{n}{X}\PY{p}{,} \PY{n}{y}\PY{p}{,} \PY{n}{C}\PY{p}{,}\PY{n}{B}\PY{o}{=}\PY{l+m+mi}{100}\PY{p}{,} \PY{n}{test\PYZus{}size}\PY{o}{=}\PY{l+m+mf}{0.3}\PY{p}{)}\PY{p}{:}
    
    \PY{n}{err} \PY{o}{=} \PY{p}{[}\PY{p}{]}
    \PY{n}{n} \PY{o}{=} \PY{n+nb}{len}\PY{p}{(}\PY{n}{y}\PY{p}{)}
    
    \PY{n}{np}\PY{o}{.}\PY{n}{random}\PY{o}{.}\PY{n}{seed}\PY{p}{(}\PY{l+m+mi}{123}\PY{p}{)}
    \PY{k}{for} \PY{n}{b} \PY{o+ow}{in} \PY{n+nb}{range}\PY{p}{(}\PY{n}{B}\PY{p}{)}\PY{p}{:}
        \PY{n}{boot} \PY{o}{=} \PY{n}{np}\PY{o}{.}\PY{n}{random}\PY{o}{.}\PY{n}{choice}\PY{p}{(}\PY{n}{n}\PY{p}{,} \PY{n}{size}\PY{o}{=}\PY{n}{n}\PY{p}{,} \PY{n}{replace}\PY{o}{=}\PY{k+kc}{True}\PY{p}{)}
        \PY{n}{X\PYZus{}b} \PY{o}{=} \PY{n}{X}\PY{p}{[}\PY{n}{boot}\PY{p}{]}
        \PY{n}{y\PYZus{}b} \PY{o}{=} \PY{n}{y}\PY{p}{[}\PY{n}{boot}\PY{p}{]}
        
        \PY{n}{X\PYZus{}train\PYZus{}b}\PY{p}{,} \PY{n}{X\PYZus{}test\PYZus{}b}\PY{p}{,} \PY{n}{y\PYZus{}train\PYZus{}b}\PY{p}{,} \PY{n}{y\PYZus{}test\PYZus{}b} \PY{o}{=} \PY{n}{train\PYZus{}test\PYZus{}split}\PY{p}{(}\PY{n}{X\PYZus{}b}\PY{p}{,} \PY{n}{y\PYZus{}b}\PY{p}{,} \PY{n}{test\PYZus{}size}\PY{o}{=}\PY{n}{test\PYZus{}size}\PY{p}{)}
        \PY{n}{model} \PY{o}{=} \PY{n}{LogisticRegression}\PY{p}{(} \PY{n}{C}\PY{o}{=}\PY{n}{C}\PY{p}{,}\PY{n}{fit\PYZus{}intercept}\PY{o}{=}\PY{k+kc}{True}\PY{p}{,} \PY{n}{random\PYZus{}state}\PY{o}{=}\PY{l+m+mi}{0}\PY{p}{,} \PY{n}{max\PYZus{}iter}\PY{o}{=}\PY{l+m+mi}{1000}\PY{p}{)}
        
        \PY{n}{\PYZus{}}\PY{p}{,} \PY{n}{test\PYZus{}err} \PY{o}{=} \PY{n}{fit\PYZus{}and\PYZus{}evaluate}\PY{p}{(}\PY{n}{model}\PY{p}{,} \PY{n}{X\PYZus{}train\PYZus{}b}\PY{p}{,} \PY{n}{y\PYZus{}train\PYZus{}b}\PY{p}{,} \PY{n}{X\PYZus{}test\PYZus{}b}\PY{p}{,} \PY{n}{y\PYZus{}test\PYZus{}b}\PY{p}{)}
        \PY{n}{err}\PY{o}{.}\PY{n}{append}\PY{p}{(}\PY{n}{test\PYZus{}err}\PY{p}{)}

    \PY{k}{return} \PY{n}{np}\PY{o}{.}\PY{n}{mean}\PY{p}{(}\PY{n}{err}\PY{p}{)}


\PY{n}{exp\PYZus{}err\PYZus{}unreg} \PY{o}{=} \PY{n}{estimate\PYZus{}expected\PYZus{}error}\PY{p}{(}\PY{n}{X\PYZus{}scaled}\PY{p}{,} \PY{n}{y}\PY{p}{,} \PY{n}{C}\PY{o}{=}\PY{n}{np}\PY{o}{.}\PY{n}{inf}\PY{p}{,}\PY{n}{B}\PY{o}{=}\PY{l+m+mi}{100}\PY{p}{)} 
\PY{n}{exp\PYZus{}err\PYZus{}wridge} \PY{o}{=} \PY{n}{estimate\PYZus{}expected\PYZus{}error}\PY{p}{(}\PY{n}{X\PYZus{}scaled}\PY{p}{,} \PY{n}{y}\PY{p}{,} \PY{n}{C}\PY{o}{=}\PY{l+m+mi}{1}\PY{p}{,}\PY{n}{B}\PY{o}{=}\PY{l+m+mi}{100}\PY{p}{)}
\PY{n}{exp\PYZus{}err\PYZus{}sridge} \PY{o}{=} \PY{n}{estimate\PYZus{}expected\PYZus{}error}\PY{p}{(}\PY{n}{X\PYZus{}scaled}\PY{p}{,} \PY{n}{y}\PY{p}{,} \PY{n}{C}\PY{o}{=}\PY{l+m+mf}{0.005}\PY{p}{,}\PY{n}{B}\PY{o}{=}\PY{l+m+mi}{100}\PY{p}{)}

\PY{n+nb}{print}\PY{p}{(}\PY{l+s+s2}{\PYZdq{}}\PY{l+s+s2}{=== Training and test errors ===}\PY{l+s+s2}{\PYZdq{}}\PY{p}{)}
\PY{n+nb}{print}\PY{p}{(}\PY{l+s+s2}{\PYZdq{}}\PY{l+s+se}{\PYZbs{}n}\PY{l+s+s2}{Unregularised Logistic Regression:}\PY{l+s+s2}{\PYZdq{}}\PY{p}{)}
\PY{n+nb}{print}\PY{p}{(}\PY{l+s+sa}{f}\PY{l+s+s2}{\PYZdq{}}\PY{l+s+s2}{Training error : }\PY{l+s+si}{\PYZob{}}\PY{n}{train\PYZus{}err\PYZus{}unreg}\PY{l+s+si}{:}\PY{l+s+s2}{.6f}\PY{l+s+si}{\PYZcb{}}\PY{l+s+s2}{\PYZdq{}}\PY{p}{)}
\PY{n+nb}{print}\PY{p}{(}\PY{l+s+sa}{f}\PY{l+s+s2}{\PYZdq{}}\PY{l+s+s2}{Test error     : }\PY{l+s+si}{\PYZob{}}\PY{n}{test\PYZus{}err\PYZus{}unreg}\PY{l+s+si}{:}\PY{l+s+s2}{.6f}\PY{l+s+si}{\PYZcb{}}\PY{l+s+s2}{\PYZdq{}}\PY{p}{)}

\PY{n+nb}{print}\PY{p}{(}\PY{l+s+s2}{\PYZdq{}}\PY{l+s+se}{\PYZbs{}n}\PY{l+s+s2}{Ridge Logistic Regression (L2) with C = 1:}\PY{l+s+s2}{\PYZdq{}}\PY{p}{)}
\PY{n+nb}{print}\PY{p}{(}\PY{l+s+sa}{f}\PY{l+s+s2}{\PYZdq{}}\PY{l+s+s2}{Training error : }\PY{l+s+si}{\PYZob{}}\PY{n}{train\PYZus{}err\PYZus{}wridge}\PY{l+s+si}{:}\PY{l+s+s2}{.6f}\PY{l+s+si}{\PYZcb{}}\PY{l+s+s2}{\PYZdq{}}\PY{p}{)}
\PY{n+nb}{print}\PY{p}{(}\PY{l+s+sa}{f}\PY{l+s+s2}{\PYZdq{}}\PY{l+s+s2}{Test error     : }\PY{l+s+si}{\PYZob{}}\PY{n}{test\PYZus{}err\PYZus{}wridge}\PY{l+s+si}{:}\PY{l+s+s2}{.6f}\PY{l+s+si}{\PYZcb{}}\PY{l+s+s2}{\PYZdq{}}\PY{p}{)}

\PY{n+nb}{print}\PY{p}{(}\PY{l+s+s2}{\PYZdq{}}\PY{l+s+se}{\PYZbs{}n}\PY{l+s+s2}{Ridge Logistic Regression (L2) with C = 0.005:}\PY{l+s+s2}{\PYZdq{}}\PY{p}{)}
\PY{n+nb}{print}\PY{p}{(}\PY{l+s+sa}{f}\PY{l+s+s2}{\PYZdq{}}\PY{l+s+s2}{Training error : }\PY{l+s+si}{\PYZob{}}\PY{n}{train\PYZus{}err\PYZus{}sridge}\PY{l+s+si}{:}\PY{l+s+s2}{.6f}\PY{l+s+si}{\PYZcb{}}\PY{l+s+s2}{\PYZdq{}}\PY{p}{)}
\PY{n+nb}{print}\PY{p}{(}\PY{l+s+sa}{f}\PY{l+s+s2}{\PYZdq{}}\PY{l+s+s2}{Test error     : }\PY{l+s+si}{\PYZob{}}\PY{n}{test\PYZus{}err\PYZus{}sridge}\PY{l+s+si}{:}\PY{l+s+s2}{.6f}\PY{l+s+si}{\PYZcb{}}\PY{l+s+s2}{\PYZdq{}}\PY{p}{)}

\PY{n+nb}{print}\PY{p}{(}\PY{l+s+s2}{\PYZdq{}}\PY{l+s+se}{\PYZbs{}n}\PY{l+s+s2}{=== Estimated Expected Generalisation Error (simulation) ===}\PY{l+s+s2}{\PYZdq{}}\PY{p}{)}
\PY{n+nb}{print}\PY{p}{(}\PY{l+s+sa}{f}\PY{l+s+s2}{\PYZdq{}}\PY{l+s+s2}{Unregularised                : }\PY{l+s+si}{\PYZob{}}\PY{n}{exp\PYZus{}err\PYZus{}unreg}\PY{l+s+si}{:}\PY{l+s+s2}{.6f}\PY{l+s+si}{\PYZcb{}}\PY{l+s+s2}{\PYZdq{}}\PY{p}{)}
\PY{n+nb}{print}\PY{p}{(}\PY{l+s+sa}{f}\PY{l+s+s2}{\PYZdq{}}\PY{l+s+s2}{Lightly Regularised (C=1)    :   }\PY{l+s+si}{\PYZob{}}\PY{n}{exp\PYZus{}err\PYZus{}wridge}\PY{l+s+si}{:}\PY{l+s+s2}{.6f}\PY{l+s+si}{\PYZcb{}}\PY{l+s+s2}{\PYZdq{}}\PY{p}{)}
\PY{n+nb}{print}\PY{p}{(}\PY{l+s+sa}{f}\PY{l+s+s2}{\PYZdq{}}\PY{l+s+s2}{Heavily Regularised (C=0.005):   }\PY{l+s+si}{\PYZob{}}\PY{n}{exp\PYZus{}err\PYZus{}wridge}\PY{l+s+si}{:}\PY{l+s+s2}{.6f}\PY{l+s+si}{\PYZcb{}}\PY{l+s+s2}{\PYZdq{}}\PY{p}{)}
\end{Verbatim}
\end{tcolorbox}

    \begin{Verbatim}[commandchars=\\\{\}]
=== Training and test errors ===

Unregularised Logistic Regression:
Training error : 0.275714
Test error     : 0.303333

Ridge Logistic Regression (L2) with C = 1:
Training error : 0.274286
Test error     : 0.303333

Ridge Logistic Regression (L2) with C = 0.005:
Training error : 0.270000
Test error     : 0.310000

=== Estimated Expected Generalisation Error (simulation) ===
Unregularised                : 0.278000
Lightly Regularised (C=1)    :   0.277667
Heavily Regularised (C=0.005):   0.277667
    \end{Verbatim}

    The training error which is a downword biased estimator for the
generalisation error, decreases as regularisation increases. The test
errors, which provides a much more accurate reflection of the
generalisation performance of a fitted model, indicate the regularised
model with hyperparameter C=0.005, generalised the worst. This implies
the model complixity was penalised too much and therefore underfits the
data. In contrast the estimated expected generalisation error of the
unregularized model, is higher than both the regularized models,
indicating the regularised leaning algorithm generalises better.

From this example, it is evident that regularisation should carefully be
conducted as too little or too much regularisation could lead to
underfitting or overfitting. To find the optimal balance the
bias-variance trade-off should be examined.

    \subsubsection{Bias-variance trade-off and the curse of high
dimensionality}\label{bias-variance-trade-off-and-the-curse-of-high-dimensionality}

Bias refers to how close an estimate is to the true value. When model
complexity is increased, there is a greater risk for overfitting and the
bias of the model is decreased. Variance refers to how stable an
estimate is. When model complexity is increased and the model starts to
overfit the data, the predictions for unseen data will be less stable
and the variance of the model will increase Hastie et al.~(2009).

It follows that when bias is decreased by increasing model complexity,
the variance will increase. This implies a negative relationship between
bias and variance, which is known as the bias-variance trade-off. When a
model has a large bias but low variance, underfitting occurs, causing
poor generalisation. Conversely, when a model has a low bias but large
variance, overfitting occurs, which also leads to poor generalisation.
The aim, therefore, in statistical learning problems, is to find a model
that minimises the generalisation error, which can only be achieved by
minimising both bias and variance.

This trade-off is found in both regression and classification. In
regression problems, the bias and variance can be minimised by
minimising the mean squared error (MSE). This error can be expressed as
\(MSE(\hat{f})=bias(\hat{f})^2+var(\hat{f})\), which is a useful
decomposition of the bias-variance trade-off. In this expression
\(bias(\hat{f})^2 = E_{\mathbf{X}}\left[(\ddot{f}(\mathbf{X}) - E_{\mathcal{D}}[\hat{f}(\mathbf{X},\mathcal{D})])^2\right]\)
and
\(var(\hat{f}) = E_{\mathbf{X}}\left[Var_{\mathcal{D}}[\hat{f}(\mathbf{X},\mathcal{D})]\right]\).

In classification, decomposing the bias-variance trade-off is not as
simple. Various metrics have been developed over the years, but none as
simple and elegant as the MSE. Domingos (2000) discovered that for
certain loss functions such as the squared error and 0-1 loss, the
generalisation error can express the bias-variance decomposition as
\(Err(\hat{f},\mathcal{D})=E_{\mathbf{X,Y}}[\mathcal{L}(\mathbf{Y},\hat{f}(\mathbf{X},\mathcal{D}))]=c_1Noise+Bias+c_2Variance\)
where \(c_1\) and \(c_2\) are coefficients depending on the loss
function.

It therefore follows that the hyperparameter C in our ongoing example,
should be chosen such that it minimises both the bias and variance. To
illustrate the how the choice of C influence the bias and variance, a
plot thereof was contructed for the ridge losgistic regression model
evaluated at varying values of C. From the plot, it is clear that where
the bias and variance intersected, there the optimal tradeoff of
minimized combined bias and variance is. This optimal point is where C
range between 0.03 and 0.06, implying when \(C\in[0.03,0.06]\),
generalisation performance is enhanced.

    \begin{tcolorbox}[breakable, size=fbox, boxrule=1pt, pad at break*=1mm,colback=cellbackground, colframe=cellborder]
\prompt{In}{incolor}{7}{\boxspacing}
\begin{Verbatim}[commandchars=\\\{\}]
\PY{k}{def}\PY{+w}{ }\PY{n+nf}{bias\PYZus{}variance\PYZus{}logistic}\PY{p}{(}\PY{n}{C\PYZus{}values}\PY{p}{,} \PY{n}{n}\PY{o}{=}\PY{l+m+mi}{200}\PY{p}{,} \PY{n}{B}\PY{o}{=}\PY{l+m+mi}{100}\PY{p}{,} \PY{n}{N}\PY{o}{=}\PY{l+m+mi}{1000}\PY{p}{)}\PY{p}{:}
    \PY{n}{bias\PYZus{}list} \PY{o}{=} \PY{p}{[}\PY{p}{]}
    \PY{n}{var\PYZus{}list} \PY{o}{=} \PY{p}{[}\PY{p}{]}
    
    \PY{n}{np}\PY{o}{.}\PY{n}{random}\PY{o}{.}\PY{n}{seed}\PY{p}{(}\PY{l+m+mi}{123}\PY{p}{)}
    \PY{n}{X}\PY{p}{,} \PY{n}{\PYZus{}} \PY{o}{=} \PY{n}{generate\PYZus{}data}\PY{p}{(}\PY{n}{N}\PY{p}{)}
    \PY{n}{X\PYZus{}design} \PY{o}{=} \PY{n}{np}\PY{o}{.}\PY{n}{c\PYZus{}}\PY{p}{[}\PY{n}{np}\PY{o}{.}\PY{n}{ones}\PY{p}{(}\PY{n}{X}\PY{o}{.}\PY{n}{shape}\PY{p}{[}\PY{l+m+mi}{0}\PY{p}{]}\PY{p}{)}\PY{p}{,} \PY{n}{X}\PY{p}{]} 
    \PY{n}{p\PYZus{}true} \PY{o}{=} \PY{n}{sigmoid}\PY{p}{(}\PY{n}{X\PYZus{}design} \PY{o}{@} \PY{n}{Beta}\PY{p}{)}
    
    \PY{k}{for} \PY{n}{C} \PY{o+ow}{in} \PY{n}{C\PYZus{}values}\PY{p}{:}
        \PY{n}{preds\PYZus{}all} \PY{o}{=} \PY{p}{[}\PY{p}{]}
        
        \PY{k}{for} \PY{n}{b} \PY{o+ow}{in} \PY{n+nb}{range}\PY{p}{(}\PY{n}{B}\PY{p}{)}\PY{p}{:}
            \PY{n}{X\PYZus{}train}\PY{p}{,} \PY{n}{y\PYZus{}train} \PY{o}{=} \PY{n}{generate\PYZus{}data}\PY{p}{(}\PY{n}{n}\PY{p}{)}
            \PY{n}{scaler} \PY{o}{=} \PY{n}{StandardScaler}\PY{p}{(}\PY{p}{)}
            \PY{n}{Xt\PYZus{}scaled} \PY{o}{=} \PY{n}{scaler}\PY{o}{.}\PY{n}{fit\PYZus{}transform}\PY{p}{(}\PY{n}{X\PYZus{}train}\PY{p}{)}  
            
            \PY{n}{model} \PY{o}{=} \PY{n}{LogisticRegression}\PY{p}{(} \PY{n}{C}\PY{o}{=}\PY{n}{C}\PY{p}{,} \PY{n}{fit\PYZus{}intercept}\PY{o}{=}\PY{k+kc}{True}\PY{p}{,} \PY{n}{random\PYZus{}state}\PY{o}{=}\PY{l+m+mi}{0}\PY{p}{,} \PY{n}{max\PYZus{}iter}\PY{o}{=}\PY{l+m+mi}{1000}\PY{p}{)}
            \PY{n}{model}\PY{o}{.}\PY{n}{fit}\PY{p}{(}\PY{n}{Xt\PYZus{}scaled}\PY{p}{,} \PY{n}{y\PYZus{}train}\PY{p}{)}

            \PY{n}{X\PYZus{}scaled} \PY{o}{=} \PY{n}{scaler}\PY{o}{.}\PY{n}{transform}\PY{p}{(}\PY{n}{X}\PY{p}{)}
            \PY{n}{p\PYZus{}hat} \PY{o}{=} \PY{n}{model}\PY{o}{.}\PY{n}{predict\PYZus{}proba}\PY{p}{(}\PY{n}{X\PYZus{}scaled}\PY{p}{)}\PY{p}{[}\PY{p}{:}\PY{p}{,} \PY{l+m+mi}{1}\PY{p}{]}
            \PY{n}{preds\PYZus{}all}\PY{o}{.}\PY{n}{append}\PY{p}{(}\PY{n}{p\PYZus{}hat}\PY{p}{)}
        
        \PY{n}{preds\PYZus{}all} \PY{o}{=} \PY{n}{np}\PY{o}{.}\PY{n}{array}\PY{p}{(}\PY{n}{preds\PYZus{}all}\PY{p}{)}  
        \PY{n}{p\PYZus{}mean} \PY{o}{=} \PY{n}{preds\PYZus{}all}\PY{o}{.}\PY{n}{mean}\PY{p}{(}\PY{n}{axis}\PY{o}{=}\PY{l+m+mi}{0}\PY{p}{)}
        
        \PY{n}{bias} \PY{o}{=} \PY{n}{np}\PY{o}{.}\PY{n}{mean}\PY{p}{(}\PY{p}{(}\PY{n}{p\PYZus{}mean} \PY{o}{\PYZhy{}} \PY{n}{p\PYZus{}true}\PY{p}{)}\PY{o}{*}\PY{o}{*}\PY{l+m+mi}{2}\PY{p}{)}
        \PY{n}{var} \PY{o}{=} \PY{n}{np}\PY{o}{.}\PY{n}{mean}\PY{p}{(}\PY{n}{np}\PY{o}{.}\PY{n}{var}\PY{p}{(}\PY{n}{preds\PYZus{}all}\PY{p}{,} \PY{n}{axis}\PY{o}{=}\PY{l+m+mi}{0}\PY{p}{)}\PY{p}{)}
        
        \PY{n}{bias\PYZus{}list}\PY{o}{.}\PY{n}{append}\PY{p}{(}\PY{n}{bias}\PY{p}{)}
        \PY{n}{var\PYZus{}list}\PY{o}{.}\PY{n}{append}\PY{p}{(}\PY{n}{var}\PY{p}{)}
    
    \PY{k}{return} \PY{n}{np}\PY{o}{.}\PY{n}{array}\PY{p}{(}\PY{n}{bias\PYZus{}list}\PY{p}{)}\PY{p}{,} \PY{n}{np}\PY{o}{.}\PY{n}{array}\PY{p}{(}\PY{n}{var\PYZus{}list}\PY{p}{)}


\PY{n}{C\PYZus{}values} \PY{o}{=} \PY{n}{np}\PY{o}{.}\PY{n}{logspace}\PY{p}{(}\PY{o}{\PYZhy{}}\PY{l+m+mi}{2}\PY{p}{,} \PY{l+m+mi}{2}\PY{p}{,} \PY{l+m+mi}{50}\PY{p}{)}  
\PY{n}{bias}\PY{p}{,} \PY{n}{variance} \PY{o}{=} \PY{n}{bias\PYZus{}variance\PYZus{}logistic}\PY{p}{(} \PY{n}{C\PYZus{}values}\PY{p}{,} \PY{n}{n}\PY{o}{=}\PY{l+m+mi}{200}\PY{p}{,} \PY{n}{B}\PY{o}{=}\PY{l+m+mi}{100}\PY{p}{,} \PY{n}{N}\PY{o}{=}\PY{l+m+mi}{1000}\PY{p}{)}

\PY{n}{plt}\PY{o}{.}\PY{n}{figure}\PY{p}{(}\PY{p}{)}
\PY{n}{plt}\PY{o}{.}\PY{n}{plot}\PY{p}{(}\PY{n}{C\PYZus{}values}\PY{p}{,} \PY{n}{bias}\PY{p}{,} \PY{n}{label}\PY{o}{=}\PY{l+s+s2}{\PYZdq{}}\PY{l+s+s2}{Bias\PYZca{}2}\PY{l+s+s2}{\PYZdq{}}\PY{p}{)}
\PY{n}{plt}\PY{o}{.}\PY{n}{plot}\PY{p}{(}\PY{n}{C\PYZus{}values}\PY{p}{,} \PY{n}{variance}\PY{p}{,} \PY{n}{label}\PY{o}{=}\PY{l+s+s2}{\PYZdq{}}\PY{l+s+s2}{Variance}\PY{l+s+s2}{\PYZdq{}}\PY{p}{)}
\PY{n}{plt}\PY{o}{.}\PY{n}{xscale}\PY{p}{(}\PY{l+s+s1}{\PYZsq{}}\PY{l+s+s1}{log}\PY{l+s+s1}{\PYZsq{}}\PY{p}{)}
\PY{n}{plt}\PY{o}{.}\PY{n}{xlabel}\PY{p}{(}\PY{l+s+s2}{\PYZdq{}}\PY{l+s+s2}{C}\PY{l+s+s2}{\PYZdq{}}\PY{p}{)}
\PY{n}{plt}\PY{o}{.}\PY{n}{title}\PY{p}{(}\PY{l+s+s2}{\PYZdq{}}\PY{l+s+s2}{Bias\PYZhy{}Variance Tradeoff (Regularised Logistic Regression)}\PY{l+s+s2}{\PYZdq{}}\PY{p}{)}
\PY{n}{plt}\PY{o}{.}\PY{n}{legend}\PY{p}{(}\PY{p}{)}
\PY{n}{plt}\PY{o}{.}\PY{n}{show}\PY{p}{(}\PY{p}{)}
\end{Verbatim}
\end{tcolorbox}

    \begin{center}
    \adjustimage{max size={0.9\linewidth}{0.9\paperheight}}{Ass1_1_files/Ass1_1_16_0.png}
    \end{center}
    { \hspace*{\fill} \\}
    
    \subsubsection{Model assessment and selection
techniques}\label{model-assessment-and-selection-techniques}

Thus far, different attributes of an optimal model have been discussed.
The model that minimises the empirical risk or alternatively, the
applicable loss function, should be selected for prediction. This model
should have the lowest bias and variance and should generalise well.
There are two techniques, the validation-set approach and
cross-validation, which can be used to select and assess such an optimal
model.

The validation-set approach is usually applied when enough data is
available. This technique splits all the data into two sets,
\(\mathcal{D}^*\) and the test set \(\mathcal{T}_e\). This is followed
by further splitting \(\mathcal{D}^*\) into the training set
\(\mathcal{T}_r\) and validation set \(\mathcal{V}\). The algorithm
continues by training all possible considered models on
\(\mathcal{T}_r\). It then selects the model which minimized the
empirical risk when applying models to \(\mathcal{V}\). Finally, the
selected model is assessed on \(\mathcal{T}_e\) with metrics such as the
generalisation error.

Cross-validation is similar to the validation-set approach, but performs
more data partitioning and is often used when there is not a lot of data
available. To perform model selection as well as model assessment, all
the data, as with the validation-set approach, should be split into two
sets, \(\mathcal{D}^*\) and the test set \(\mathcal{T}_e\). The
algorithm then divides \(\mathcal{D}^*\) into K mutually exclusive and
exhaustive groups \(\mathcal{V}^{(k)}:k=1,..,K\) called folds. When K=1,
this technique is called leave-one-out cross-validation (LOOCV). Next,
each possible model is trained K times, each time, k, on
\(\mathcal{D}^*\) excluding fold k. It then selects the model,
\(\hat{f}\), which minimises the cross validation error, cver
\((\hat{f},\mathcal{D},K)=\frac{1}{K}\sum^K_{k=1}\hat{\mathcal{R}}(\hat{f}(.,\mathcal{D}^*\backslash \mathcal{V}^{(k)}),\mathcal{V}^{(k)})\)
where \(\mathcal{D}^*\backslash \mathcal{V}^{(k)}\) refers
\(\mathcal{D}^*\) excluding fold k. Exactly like the validation-set
approach, the selected model is assessed on \(\mathcal{T}_e\) with
metrics such as the generalisation error.

Selecting the best model in the regularised logistic regression
application, is equavalent to finding the optimal C value. Using 10 fold
cross-validation, models with \(C\in[0.04,0.07]\), similar than before,
minimises the cross-validation error and can be regarded as sensible
model selections. Assessing these models with the test errors indicates
however that models with C values slighly smaller (\(C\in[0.02,0.05]\))
yields better generalisation performance for the currect dataset.

    \begin{tcolorbox}[breakable, size=fbox, boxrule=1pt, pad at break*=1mm,colback=cellbackground, colframe=cellborder]
\prompt{In}{incolor}{8}{\boxspacing}
\begin{Verbatim}[commandchars=\\\{\}]
\PY{k+kn}{from}\PY{+w}{ }\PY{n+nn}{sklearn}\PY{n+nn}{.}\PY{n+nn}{model\PYZus{}selection}\PY{+w}{ }\PY{k+kn}{import} \PY{n}{KFold}

\PY{n}{np}\PY{o}{.}\PY{n}{random}\PY{o}{.}\PY{n}{seed}\PY{p}{(}\PY{l+m+mi}{123}\PY{p}{)}

\PY{n}{X\PYZus{}train}\PY{p}{,} \PY{n}{X\PYZus{}test}\PY{p}{,} \PY{n}{y\PYZus{}train}\PY{p}{,} \PY{n}{y\PYZus{}test} \PY{o}{=} \PY{n}{train\PYZus{}test\PYZus{}split}\PY{p}{(}\PY{n}{X\PYZus{}scaled}\PY{p}{,} \PY{n}{y}\PY{p}{,} \PY{n}{test\PYZus{}size}\PY{o}{=}\PY{l+m+mf}{0.5}\PY{p}{,} \PY{n}{random\PYZus{}state}\PY{o}{=}\PY{l+m+mi}{0}\PY{p}{)}

\PY{k}{def}\PY{+w}{ }\PY{n+nf}{kfold\PYZus{}cv\PYZus{}error}\PY{p}{(}\PY{n}{X}\PY{p}{,} \PY{n}{y}\PY{p}{,} \PY{n}{C}\PY{p}{,} \PY{n}{K}\PY{o}{=}\PY{l+m+mi}{5}\PY{p}{)}\PY{p}{:}
    \PY{n}{kf} \PY{o}{=} \PY{n}{KFold}\PY{p}{(}\PY{n}{n\PYZus{}splits}\PY{o}{=}\PY{n}{K}\PY{p}{,} \PY{n}{shuffle}\PY{o}{=}\PY{k+kc}{True}\PY{p}{,} \PY{n}{random\PYZus{}state}\PY{o}{=}\PY{l+m+mi}{0}\PY{p}{)}
    
    \PY{n}{errors} \PY{o}{=} \PY{p}{[}\PY{p}{]}
    
    \PY{k}{for} \PY{n}{train\PYZus{}idx}\PY{p}{,} \PY{n}{val\PYZus{}idx} \PY{o+ow}{in} \PY{n}{kf}\PY{o}{.}\PY{n}{split}\PY{p}{(}\PY{n}{X}\PY{p}{)}\PY{p}{:}
        \PY{n}{X\PYZus{}tr}\PY{p}{,} \PY{n}{X\PYZus{}val} \PY{o}{=} \PY{n}{X}\PY{p}{[}\PY{n}{train\PYZus{}idx}\PY{p}{]}\PY{p}{,} \PY{n}{X}\PY{p}{[}\PY{n}{val\PYZus{}idx}\PY{p}{]}
        \PY{n}{y\PYZus{}tr}\PY{p}{,} \PY{n}{y\PYZus{}val} \PY{o}{=} \PY{n}{y}\PY{p}{[}\PY{n}{train\PYZus{}idx}\PY{p}{]}\PY{p}{,} \PY{n}{y}\PY{p}{[}\PY{n}{val\PYZus{}idx}\PY{p}{]}
        
        \PY{n}{model} \PY{o}{=} \PY{n}{LogisticRegression}\PY{p}{(} \PY{n}{C}\PY{o}{=}\PY{n}{C}\PY{p}{,} \PY{n}{fit\PYZus{}intercept}\PY{o}{=}\PY{k+kc}{True}\PY{p}{,} \PY{n}{random\PYZus{}state}\PY{o}{=}\PY{l+m+mi}{0}\PY{p}{,} \PY{n}{max\PYZus{}iter}\PY{o}{=}\PY{l+m+mi}{1000}\PY{p}{)}
        \PY{n}{model}\PY{o}{.}\PY{n}{fit}\PY{p}{(}\PY{n}{X\PYZus{}tr}\PY{p}{,} \PY{n}{y\PYZus{}tr}\PY{p}{)}
        \PY{n}{errors}\PY{o}{.}\PY{n}{append}\PY{p}{(}\PY{n}{zero\PYZus{}one\PYZus{}loss}\PY{p}{(}\PY{n}{y\PYZus{}val}\PY{p}{,} \PY{n}{model}\PY{o}{.}\PY{n}{predict}\PY{p}{(}\PY{n}{X\PYZus{}val}\PY{p}{)}\PY{p}{)}\PY{p}{)}
    
    \PY{k}{return} \PY{n}{np}\PY{o}{.}\PY{n}{mean}\PY{p}{(}\PY{n}{errors}\PY{p}{)}

\PY{n}{C\PYZus{}values} \PY{o}{=} \PY{n}{np}\PY{o}{.}\PY{n}{logspace}\PY{p}{(}\PY{o}{\PYZhy{}}\PY{l+m+mi}{2}\PY{p}{,} \PY{l+m+mi}{2}\PY{p}{,} \PY{l+m+mi}{100}\PY{p}{)}
\PY{n}{cv\PYZus{}errors} \PY{o}{=} \PY{p}{[}\PY{p}{]}
\PY{n}{test\PYZus{}errors} \PY{o}{=} \PY{p}{[}\PY{p}{]}
\PY{n}{exp\PYZus{}errors} \PY{o}{=} \PY{p}{[}\PY{p}{]}

\PY{k}{for} \PY{n}{C} \PY{o+ow}{in} \PY{n}{C\PYZus{}values}\PY{p}{:}
    \PY{n}{cv\PYZus{}errors}\PY{o}{.}\PY{n}{append}\PY{p}{(}\PY{n}{kfold\PYZus{}cv\PYZus{}error}\PY{p}{(}\PY{n}{X\PYZus{}train}\PY{p}{,} \PY{n}{y\PYZus{}train}\PY{p}{,} \PY{n}{C}\PY{p}{,} \PY{n}{K}\PY{o}{=}\PY{l+m+mi}{10}\PY{p}{)}\PY{p}{)}
    
    \PY{n}{model} \PY{o}{=} \PY{n}{LogisticRegression}\PY{p}{(} \PY{n}{C}\PY{o}{=}\PY{n}{C}\PY{p}{,} \PY{n}{fit\PYZus{}intercept}\PY{o}{=}\PY{k+kc}{True}\PY{p}{,} \PY{n}{random\PYZus{}state}\PY{o}{=}\PY{l+m+mi}{0}\PY{p}{,} \PY{n}{max\PYZus{}iter}\PY{o}{=}\PY{l+m+mi}{1000}\PY{p}{)}
    \PY{n}{model}\PY{o}{.}\PY{n}{fit}\PY{p}{(}\PY{n}{X\PYZus{}train}\PY{p}{,} \PY{n}{y\PYZus{}train}\PY{p}{)}
    \PY{n}{test\PYZus{}errors}\PY{o}{.}\PY{n}{append}\PY{p}{(}\PY{n}{zero\PYZus{}one\PYZus{}loss}\PY{p}{(}\PY{n}{y\PYZus{}test}\PY{p}{,} \PY{n}{model}\PY{o}{.}\PY{n}{predict}\PY{p}{(}\PY{n}{X\PYZus{}test}\PY{p}{)}\PY{p}{)}\PY{p}{)}

    \PY{n}{exp\PYZus{}errors}\PY{o}{.}\PY{n}{append}\PY{p}{(}\PY{n}{estimate\PYZus{}expected\PYZus{}error}\PY{p}{(}\PY{n}{X\PYZus{}scaled}\PY{p}{,} \PY{n}{y}\PY{p}{,} \PY{n}{C}\PY{o}{=}\PY{n}{C}\PY{p}{,}\PY{n}{B}\PY{o}{=}\PY{l+m+mi}{100}\PY{p}{)}\PY{p}{)}


\PY{n}{plt}\PY{o}{.}\PY{n}{figure}\PY{p}{(}\PY{p}{)}
\PY{n}{plt}\PY{o}{.}\PY{n}{plot}\PY{p}{(}\PY{n}{C\PYZus{}values}\PY{p}{,} \PY{n}{cv\PYZus{}errors}\PY{p}{,} \PY{n}{label}\PY{o}{=}\PY{l+s+s2}{\PYZdq{}}\PY{l+s+s2}{CV Error}\PY{l+s+s2}{\PYZdq{}}\PY{p}{)}
\PY{n}{plt}\PY{o}{.}\PY{n}{plot}\PY{p}{(}\PY{n}{C\PYZus{}values}\PY{p}{,} \PY{n}{test\PYZus{}errors}\PY{p}{,} \PY{n}{label}\PY{o}{=}\PY{l+s+s2}{\PYZdq{}}\PY{l+s+s2}{Test Error}\PY{l+s+s2}{\PYZdq{}}\PY{p}{)}
\PY{n}{plt}\PY{o}{.}\PY{n}{xscale}\PY{p}{(}\PY{l+s+s1}{\PYZsq{}}\PY{l+s+s1}{log}\PY{l+s+s1}{\PYZsq{}}\PY{p}{)}
\PY{n}{plt}\PY{o}{.}\PY{n}{xlabel}\PY{p}{(}\PY{l+s+s2}{\PYZdq{}}\PY{l+s+s2}{C}\PY{l+s+s2}{\PYZdq{}}\PY{p}{)}
\PY{n}{plt}\PY{o}{.}\PY{n}{ylabel}\PY{p}{(}\PY{l+s+s2}{\PYZdq{}}\PY{l+s+s2}{Error}\PY{l+s+s2}{\PYZdq{}}\PY{p}{)}
\PY{n}{plt}\PY{o}{.}\PY{n}{title}\PY{p}{(}\PY{l+s+s2}{\PYZdq{}}\PY{l+s+s2}{Cross\PYZhy{}Validation vs Test Error}\PY{l+s+s2}{\PYZdq{}}\PY{p}{)}
\PY{n}{plt}\PY{o}{.}\PY{n}{legend}\PY{p}{(}\PY{p}{)}
\PY{n}{plt}\PY{o}{.}\PY{n}{show}\PY{p}{(}\PY{p}{)}
\end{Verbatim}
\end{tcolorbox}

    \begin{center}
    \adjustimage{max size={0.9\linewidth}{0.9\paperheight}}{Ass1_1_files/Ass1_1_18_0.png}
    \end{center}
    { \hspace*{\fill} \\}
    
    Interestingly, the cross-validation error can also serve as an estimator
for the expected generalisation error. This is because models are
trained on K different datasets and evaluated on unseen data. Again, the
bias-variance trade-off can be observed when choosing the value of K.
When K is large, the bias is small and the variance is large since
models will be trained and evaluated on more diverse datasets. This
concept is visualised with the following plot.

    \begin{tcolorbox}[breakable, size=fbox, boxrule=1pt, pad at break*=1mm,colback=cellbackground, colframe=cellborder]
\prompt{In}{incolor}{9}{\boxspacing}
\begin{Verbatim}[commandchars=\\\{\}]
\PY{n}{plt}\PY{o}{.}\PY{n}{figure}\PY{p}{(}\PY{p}{)}
\PY{n}{plt}\PY{o}{.}\PY{n}{plot}\PY{p}{(}\PY{n}{C\PYZus{}values}\PY{p}{,} \PY{n}{cv\PYZus{}errors}\PY{p}{,} \PY{n}{label}\PY{o}{=}\PY{l+s+s2}{\PYZdq{}}\PY{l+s+s2}{CV Error}\PY{l+s+s2}{\PYZdq{}}\PY{p}{)}
\PY{n}{plt}\PY{o}{.}\PY{n}{plot}\PY{p}{(}\PY{n}{C\PYZus{}values}\PY{p}{,} \PY{n}{exp\PYZus{}errors}\PY{p}{,} \PY{n}{label}\PY{o}{=}\PY{l+s+s2}{\PYZdq{}}\PY{l+s+s2}{Estimated Expected Generalisation Error}\PY{l+s+s2}{\PYZdq{}}\PY{p}{)}
\PY{n}{plt}\PY{o}{.}\PY{n}{xscale}\PY{p}{(}\PY{l+s+s1}{\PYZsq{}}\PY{l+s+s1}{log}\PY{l+s+s1}{\PYZsq{}}\PY{p}{)}
\PY{n}{plt}\PY{o}{.}\PY{n}{xlabel}\PY{p}{(}\PY{l+s+s2}{\PYZdq{}}\PY{l+s+s2}{C}\PY{l+s+s2}{\PYZdq{}}\PY{p}{)}
\PY{n}{plt}\PY{o}{.}\PY{n}{ylabel}\PY{p}{(}\PY{l+s+s2}{\PYZdq{}}\PY{l+s+s2}{Error}\PY{l+s+s2}{\PYZdq{}}\PY{p}{)}
\PY{n}{plt}\PY{o}{.}\PY{n}{title}\PY{p}{(}\PY{l+s+s2}{\PYZdq{}}\PY{l+s+s2}{Cross\PYZhy{}Validation vs Test Error}\PY{l+s+s2}{\PYZdq{}}\PY{p}{)}
\PY{n}{plt}\PY{o}{.}\PY{n}{legend}\PY{p}{(}\PY{p}{)}
\PY{n}{plt}\PY{o}{.}\PY{n}{show}\PY{p}{(}\PY{p}{)}
\end{Verbatim}
\end{tcolorbox}

    \begin{center}
    \adjustimage{max size={0.9\linewidth}{0.9\paperheight}}{Ass1_1_files/Ass1_1_20_0.png}
    \end{center}
    { \hspace*{\fill} \\}
    
    \subsubsection{Conclusion}\label{conclusion}

Supervised learning provides a structured framework for predicting
responses from observed features, with the choice of loss functions,
regularisation, and model selection techniques playing critical roles in
achieving optimal performance. Through the logistic regression example,
it was shown that while regularisation can improve generalisation, too
much penalisation leads to underfitting, and too little risks
overfitting. Careful evaluation using training, test, and expected
generalisation errors, alongside bias-variance analysis and
cross-validation, ensures models achieve the right balance between bias
and variance. Ultimately, the effectiveness of supervised learning lies
in selecting models that minimise generalisation error while maintaining
robustness and interpretability.

    


    % Add a bibliography block to the postdoc
    
    
    
\end{document}
